%% Leakage Horizons in Group Membership Prediction.
%% Generated from paper/penumbra-revised.docx by scripts/build_jmlr.py -- edit the source, not this.
\documentclass[twoside,11pt]{article}

\usepackage{jmlr2e}

%% jmlr2e loads epsfig, amssymb, natbib, graphicx and hyperref. These are the ones it does not.
\usepackage[T1]{fontenc}     % accented letters as glyphs, and hyphenation across them
\usepackage{amsmath}          % \underbrace with limits, \substack, \binom
\usepackage{booktabs}         % \toprule, \midrule, \bottomrule in the tables
\usepackage{longtable}        % the tables that cannot fit a page stay breakable
\usepackage{array}            % >{\raggedright\arraybackslash} column specs
\usepackage{calc}             % \real{} in the column widths

%% Pandoc emits these and defines them in its own preamble, which this document does not use.
\providecommand{\tightlist}{%
  \setlength{\itemsep}{0pt}\setlength{\parskip}{0pt}}

%% Heading arguments are {volume}{year}{pages}{date submitted}{date published}{paper id}{authors}
\jmlrheading{1}{2026}{1--48}{8/26}{}{26-0000}{M\"ucahit Kutsal, Alpaslan Erda\u{g}, U\u{g}ur \"Ozbek and Do\u{g}a Ko\c{c}yi\u{g}it}

%% The running head is set in small caps, and in these fonts the small capital of "i" has no dot. In
%% English that is correct. In Turkish the dotted and the dotless i are different letters, so a name
%% like Ko\c{c}yi\u{g}it comes out misspelled there. The class's own documented short author form
%% is "Minton et al", which is what this is, and the full list keeps its dots in the footer, which is
%% not small caps.
\ShortHeadings{Leakage Horizons in Group Membership Prediction}{Kutsal et al.}
\firstpageno{1}

\begin{document}

\title{Leakage Horizons in Group Membership Prediction}

\author{\name M\"ucahit Kutsal \email mucahitkutsal@ankara.edu.tr \\
       \addr Ankara University, Ankara, T\"urkiye \\
       Republic of T\"urkiye Ministry of Foreign Affairs, Ankara, T\"urkiye
       \AND
       \name Alpaslan Erda\u{g} \email alpaslan.erdag@ankara.edu.tr \\
       \addr Ankara University, Ankara, T\"urkiye \\
       Republic of T\"urkiye Ministry of Foreign Affairs, Ankara, T\"urkiye
       \AND
       \name U\u{g}ur \"Ozbek \email uozbek@mfa.gov.tr \\
       \addr Republic of T\"urkiye Ministry of Foreign Affairs, Ankara, T\"urkiye \\
       Konya Technical University, Konya, T\"urkiye
       \AND
       \name Do\u{g}a Ko\c{c}yi\u{g}it \email doga.kocyigit@stud.tu-darmstadt.de \\
       \addr Technical University of Darmstadt, Darmstadt, Germany}

\editor{}

\maketitle

\begin{abstract}%
Group membership forecasters draw relational features from the registers that produced their labels, so a register dating an event early encodes the label before it exists. A correct temporal split does not remove this: the contamination sits inside the register, not across it. We formalise its reach as a \textbf{leakage horizon}, the smallest lead time at which the target relation carries no administrative trace of the event, and build diagnostics to estimate it. Against a width a generator sets, and therefore knows, the estimate lands on it in all six settings and never short of it: a conservative estimate, not the lower bound we specified, whose upper end is not identified. The suite also misreports. A statistic pooled over the whole target relation cannot tell a trace of the event from a standing similarity between the joining entity and the members: concentrating that similarity drives its slope on the true width from 1.00 to 0.00, and widening the range does not bring it back. A statistic read on the group's own record stays unbiased throughout. In a trade register that codes both a signature and an entry-into-force date, the back-dated dyads are 6.8 percent of the relation, invisible to the pooled statistic. On a coalition register the pooled reading still shows a contrast after eight years, where the group's own record carries nothing past one year. Eleven ways such a diagnostic misreports are listed and ten demonstrated, four of them on a protein complex benchmark.
\end{abstract}

\begin{keywords}
label leakage; temporal validation; group membership prediction; graph neural networks; provenance; measurement validity
\end{keywords}

\section{Introduction}

A group membership forecaster takes a set of candidate entities, a group, and a body of relational data, and asks whether the candidate will join. The pattern recurs across domains. Recommender systems predict community joins from social graphs. Computational biology predicts protein complex membership from interaction networks. Political forecasting predicts accession to international organisations from trade, alliance and voting networks. In every case the relational data and the labels come from curated registers, and in every case those registers are compiled by people who record events after the fact.

This creates a failure mode that ordinary temporal validation does not catch. Suppose a register records an organisation's membership by year, and suppose that for some organisations it lists members from the year a founding convention was signed rather than the year the organisation began operating. A system trained on relational features derived from that register, and evaluated with a strict temporal split, will still see the membership edge before the coded accession date. The split is honest. The features are contaminated anyway, because the contamination is inside the register rather than across the split.

The quantity that matters is not whether such contamination exists but how far it reaches. A system forecasting one year ahead may be entirely compromised while the same system forecasting five years ahead is not. We call the smallest lead time at which the target relation no longer carries administrative traces of the event the \textbf{leakage horizon}, denoted \(\Delta^{\star}\). Estimating it is a measurement problem, and this paper is about the instrument.

Our contribution is in five parts.

\textbf{A diagnostic suite, and one conservative estimate rather than an interval.} Section 4 defines a suite of four components over a multi-relational temporal graph, three of which follow the usual leakage-audit design and one of which, focal isolation, we show the design cannot do without. The horizon was specified as an interval, bracketed below by an exact clique statistic and above by rank position and single-feature discrimination, on the expectation that the two families have opposite error profiles: specific and insensitive against sensitive and non-specific. Section 6.1 reports that the bracket does not hold. The upper end as written is unreachable in half the synthetic settings and never fires on the field corpus; rewritten as a band crossing it fires everywhere, coincides with the lower end, and on the field corpus falls below it. What survives is a single number whose errors run one way: it never undershoots the width the generator set and overshoots by one or two periods in a third of resamples at the widest setting. That is not the profile of a lower bound, since a lower bound above the quantity it bounds has failed rather than been careful, so Section 3.2 states what it is instead, a conservative estimate under a named loss.

\textbf{Calibration against two independent ground truths.} Section 6.1 calibrates the estimate against a back-dating parameter set by a synthetic generator whose relational density is held exactly invariant across settings, and recovers it in 6 of 6 settings, with a resampling distribution that never falls below the target. Section 6.4 calibrates against a real register that codes two dates per agreement, signature and entry into force, so the back-dating width is printed in the source rather than estimated, over 759 agreements. The instrument works in the first case; in the second the back-dated dyads are 6.8 percent of the relation and the pooled statistic the suite specifies is dominated by the other 93.2 percent. Those two are where the pooled reading is falsified against a width that is known, and they are the only two. Section 6.5 shows that no stratum of the coalition corpus supports both instruments as measurements at once, so what that corpus establishes is a different and weaker thing, set out where it is reported. The gap between those two outcomes is the paper's central finding, and it is a gap between two instruments read on one set of events rather than a ratio between two quantities in one unit.

\textbf{A measured identification failure, not an argued one.} Section 6.2 gives the reason for that gap as a curve rather than as reasoning. On a generator whose homophily strength is a controlled parameter and whose back-dating width is set, the pooled statistic's estimate loses all dependence on the true width as homophily rises. The statistic computed on the focal group's own record stays exactly unbiased at every setting. The collapse is not an artefact of the range examined. The censoring rate is reported per cell and every unit slope in the table is uncensored. Rerunning the collapsed cells over a range two and a half times as wide moves the estimate with the range without returning its dependence on the truth, which goes from zero to about minus one. Nor is it a property of one functional form, though establishing that took a second measurement. Homophily is a power law on latent distance in the headline grid, and we reran the whole grid under a bounded kernel and under discrete latent blocks. The collapse appears in two of the three and not in the third, which says nothing about the kernels and everything about the axis: each scales the parameter differently, and the kernel that does not collapse never reaches the concentration at which the others do. Reported instead against a quantity a field corpus also has, the two kernels that reach it break the statistic at anchor clique contrasts of 0.0152 and 0.0186, and corpus A's own strongest anchor is 0.0259. The same experiment shows that the anchor band repairs the pooled statistic only when the anchors share the target relation's latent structure, and that saturation is the condition under which homophily breaks the clique statistic, which joins two of the mechanisms below into one.

\textbf{A condition on the operation we recommend.} Focal isolation reads one group's record, and the case for it is that this is the only place administrative back-dating can live. Section 6.3 puts that claim on the same generator with a curatorial wave as a controlled parameter. It holds under one dating convention and fails under mixed ones: where a neighbouring group's register is coarser than the focal group's, the trace reaches into the target relation from a record the focal reading never opens. The remedy is procedural rather than statistical and it is measured, not proposed: reading every register that feeds the target relation, rather than the focal group's alone, closes most of the shortfall.

\textbf{Ten failure modes demonstrated, and one we could not.} Section 7 states each mechanism by which the diagnostic misreports, with the measurement that establishes it and the kind of claim it is. Eleven are listed and ten are demonstrated: the eleventh is a concern we tested and could not demonstrate, and Section 7 sets it apart rather than counting it. Because several are claims about any curated group corpus, Section 6.6 tests them on a second domain chosen for having the inverse group size distribution. One claim is falsified there and restated, one is narrowed, and one has its explanation reversed by the controlled generator. This matters beyond our diagnostic: several of the mechanisms are properties of the underlying statistics rather than of our implementation, and they apply to any system that measures relational leakage by comparing feature densities between positive and negative groups.

Section 9 reports the audit that produced the largest number of corrections, and it is not about leakage at all. Nine defects were found in a single working day by asking four questions with externally checkable answers, none of which had been caught by a validation suite that agreed with itself, two independent count audits, or a checker that compared every document number against its source artefact. We state the four questions, because they transfer to any pipeline of this shape, together with a fifth and a sixth that preparing this revision showed the four do not contain, the second of which cannot be automated at all.

\section{Related Work}

\textbf{Label leakage and target leakage.} The general problem of features that encode the label is long recognised in applied machine learning, and \citet{kaufman} give the canonical treatment, distinguishing leakage in training examples from leakage in feature construction. Their prescription is a separation between the data available at prediction time and the data used to build features. The concrete failures of that separation were already being catalogued in the cross-validation literature, where \citet{rao} show how a resampling scheme that looks correct can leak through preprocessing. That prescription is necessary and, in the setting we study, insufficient: the contamination is not a violation of the temporal boundary but a property of how the register dates its own contents, so a correctly implemented boundary does not remove it. \citet{kapoor_leakage} survey 294 papers across seventeen fields and give a taxonomy with eight leaves, and where our case sits in it is worth stating precisely, because the two obvious guesses are both wrong. It is not under L1, a lack of clean separation between training and test data, whose four leaves are an absent test set, preprocessing across the split, feature selection across the split and duplicated records; none of them is present here. The horizon belongs under \textbf{L2, a model using features that are not legitimate}, and it is a subtype L2 does not name. In their examples the illegitimacy is the analyst's: a feature chosen that stands proxy for the target. Here the analyst's choice is defensible and the illegitimacy is the register's, because a feature computed from today's register at a historical origin is not a feature that existed at that origin. It is also worth saying what it is not, since the obvious guess is wrong: it is not L3.1, temporal leakage, whose definition is a split that is not chronological. Ours is chronological and the contamination survives it, which is the whole point. The two remaining leaves are equally not it. L3.2 is dependence between training and test samples, which our resampling unit addresses separately and which would not create a decaying horizon curve. L3.3 is a test distribution sampled unrepresentatively, which is a question about the comparison stratum rather than about dating. All eight were checked and this is the result. Their remedy is a model info sheet documenting how the split was constructed. Our declaration protocol is that idea pushed one step further: what is documented in advance is not only the split but what each relation records, which is what would make it a prospective test rather than a description. We report below how far short of that we fell for one of the two corpora, because the failure is in the documenting rather than in the idea. Their consensus reporting standard, REFORMS \citep{reforms}, is the natural home for the questions Section 9 reports.

\textbf{Temporal validation of graph learning.} Work on evaluation protocols for graph neural networks has repeatedly found that reported gains shrink or vanish under stricter protocols. \citet{shchur} show that architectural comparisons on citation benchmarks are dominated by splitting and tuning choices. \citet{ogb} introduce benchmarks with explicit temporal splits for exactly this reason, and \citet{tgb} extend that to temporal graphs specifically, with negative sampling and evaluation protocols designed around the fact that a temporal split alone does not make a benchmark hard. \citet{yang_lp} make the same point one level down, for link prediction specifically: the class imbalance and the sampling of the negative set decide the reported metric more than the method does. That is why we state corpus B's negative construction in full and make the comparison stratum part of the estimand. Our finding is complementary and sharper in one respect: we exhibit a corpus where a strict temporal split is in place, the split is correct, and a leakage diagnostic applied on top of it still reports an eight-year contamination that does not exist.

\textbf{Recurring edges, and the priority of mechanism (i).} \citet{poursafaei} show that a large share of the edges in standard dynamic link prediction benchmarks are repetitions of edges seen earlier, and that a memorisation baseline which predicts previously observed edges is competitive with trained models on many of them. They also show that the negative sampling strategy decides much of the reported performance. Our mechanism (i) is the group-membership form of the first of those observations. An entity that left a group and rejoined is on the register at long lead times for a reason that is not anticipation, and reaching further back is more likely to land inside an earlier spell. We claim no priority for the observation that repetition inflates temporal evaluation. What we add is that in a \emph{leakage diagnostic} the confound does not merely inflate a score: it manufactures a decaying curve with the shape of a horizon. It also loads onto whichever stratum has the higher recurrence rate, which in our corpus is also the stratum whose dates are least precise. Their negative sampling discussion bears directly on corpus B, whose strata decide the level of every contrast we report.

\textbf{Measurement validity.} The paper is a construct validity argument throughout: it asks whether a statistic measures the quantity its definition names. The vocabulary is \citet{cronbach}'s, who introduced construct validity as the case where a test has no criterion to be checked against and its meaning has to be established by the pattern of relations it enters. That is the position a leakage diagnostic is in: there is no gold-standard horizon to correlate with. \citet{jacobs} carry that framework into machine learning, distinguishing the unobservable construct from the operationalisation that stands in for it and enumerating the ways the substitution fails. Our separation of what is measured from what is entailed is that distinction applied to one estimand. Their vocabulary is also what makes our central field claim statable: a pooled contrast over the target relation and the leakage horizon are not the same construct, so the identification failure is a convergent validity failure rather than a bug. Two further results in that line bear on the design rather than on the interpretation. \citet{geiger} find that social-computing papers routinely do not report where their labels came from, which is the gap our declaration is meant to close for relations rather than for labels. The convergent-validity literature also has a standard move we borrow. Where two instruments claim to measure one construct and disagree, the disagreement is evidence about the instruments, and the tie is broken by an external criterion. Ours is a trade register that codes the width twice.

\textbf{Shortcut learning.} Component U2 is named a shortcut probe and mechanisms (f) and (g) are shortcut arguments, so the term should carry its source. \citet{geirhos} set out the pattern: a model that solves a benchmark through a feature correlated with the label but not with the construct. They add that shortcut reliance is a property of the training distribution as much as of the architecture, which is what mechanism (g) reports in a form we did not expect. The sharpest instances are in natural language inference, where \citet{gururangan} show that annotation artefacts let a hypothesis-only model do most of the task and \citet{mccoy} build a diagnostic set that separates a heuristic from the construct. Both are the move our U2 unit test makes, and both are instructive about its limit: a diagnostic set is available when a researcher can write the counterexamples, and a leakage horizon has no counterexamples to write, which is why our version needs a generator.

\textbf{Provenance as evidence.} A parallel literature in data management treats provenance as a first class artefact rather than as documentation. \citet{simmhan} survey its uses in e-science and \citet{cheney} give the why, how and where distinction that our provenance stamping implements in its crudest form: our artefacts record the command that produced them, which is \emph{how} provenance, and the corpus fingerprint, which is \emph{where}. The documentation standards are the other half of the literature, and they converge on recording origin as a condition of release: datasheets for datasets \citep{datasheets}, model cards for models \citep{modelcards}, and the survey by \citet{paullada} of what dataset practice actually does. \citet{sambasivan} give the cost of not doing it, in the form of data cascades that surface far downstream of the decision that caused them, which is the shape of every defect our own audit reports. Our use is narrower than any of these and, we think, underused in evaluation practice: we treat a written declaration of what each relational channel records as a \emph{testable prediction} about which channels will show leakage, and we report the test and its limit below.

\textbf{Seed variance and the reporting of margins.} Our architectural comparison turns on a point that is well established and still routinely skipped: a margin between two models is not interpretable without the spread of the seeds behind it. \citet{henderson} show for reinforcement learning that reported differences of the size usually published are inside the seed noise, and \citet{dodge} give the reporting practice that follows, which is to publish the distribution rather than the best or the mean of a small sample. Our Table 10 is an instance of their argument arriving at a different field, and the fourth of our audit questions is the same argument turned into an audit step.

\textbf{Group and coalition prediction.} Hyperedge prediction and higher-order link prediction supply the task formulation. \citet{benson} give the simplicial closure formulation and show that higher-order structure predicts higher-order events better than its pairwise projection does. \citet{do_hyper} characterise the structural regularities real hypergraphs have that random ones do not, and \citet{yoon} quantify how much of the predictive gain survives when the comparison is made carefully, which is the same question our capacity-matched controls ask of message passing. The encoders we compare are standard: neighbourhood aggregation in the sense of \citet{hamilton}, its relation-specific form from \citet{rgcn}, and a temporal encoder over the graph history in the sense of \citet{tgn}. The specific instance we use for the field application is accession to intergovernmental organisations, for which the Correlates of War project supplies a membership register with documented coding rules \citep{cow_igo}, descended from the version 2 release \citep{cow_igo2} whose documentation states the dating conventions this paper spends its length measuring. The substantive literature on why such organisations end, which is what Appendix A's termination task estimates, is small and \citet{eilstrup} is its most systematic entry. The protein complex testbed uses CORUM \citep{corum} and Complex Portal \citep{complexportal} complexes against the STRING interaction network \citep{string}, a standard pairing in complex prediction. The community assessment of module identification on such networks \citep{choobdar} is the reason we treat the curated-register channel as the label source rather than as one predictor among ten.

\textbf{Competing risks and interval censoring.} Our termination task is a discrete-time competing-risks problem with interval-censored event times, and we use a likelihood contribution of the form \(S(t_{L}) - S(t_{R})\) over entity periods. Each of those three words carries a standard treatment and the combination is what the register forces. Proportional hazards is \citet{cox}'s. The discrete-time formulation we actually estimate, a sequence of conditional probabilities over entity periods, is \citet{allison}'s. Competing risks in the cause-specific form we use are \citet{prentice}'s; \citet{finegray} give the subdistribution alternative, which we do not use because our question is about cause-specific mechanisms rather than about cumulative incidence. Interval censoring in a proportional hazards model is \citet{finkelstein}'s. \citet{boxsteff} is the standard reference for the same machinery in political science, which is the substantive literature Appendix A's coefficients speak to. We note the combination here because the register's own coding rule, which declares an organisation defunct after ten consecutive years without a plenary session, makes the interval formulation mandatory rather than optional: the event is known to have happened inside a window and not at a date.

\section{Problem Formulation}

This section fixes the objects the diagnostics act on: a corpus and the dating strata its register imposes, the horizon itself, the two operations the suite applies to a focal group, and the line between what those operations measure and what they can only entail.

\subsection{Corpus and Notation}

A \textbf{group corpus} is a set of entities \(c \in C\), each with a membership set \(S_{c}(t) \subseteq V\) that varies over discrete time \(t\), together with a set of \textbf{events}. An accession event is a triple \((c,v,t)\) recording that \(v\) joined \(c\) at \(t\). Because registers record dates at varying resolution, each event carries an interval \((t_{L},t_{R}\rbrack\) and a \textbf{dating stratum}; an exactly dated event is the special case \(t_{R} = t_{L} + 1\).

An \textbf{observation graph} \(G_{t} = (V,E_{t}^{r_{1}},\ldots,E_{t}^{r_{m}})\) carries \(m\) relations over the same node set. One relation is designated \(r^{\star}\), the \textbf{target relation}, defined as the relation through which the event would leak if it leaks at all. For accession to an organisation, \(r^{\star}\) is co-membership: two entities share an edge when they belong to a common organisation in that year. Note that \(r^{\star}\) is built from the same register that supplies the labels. This is deliberate. It leaks by construction, and the question is not whether but how far back.

For a group \(S\) and relation \(r\), write \(\Pi(S)\) for the set of unordered pairs of members and define the \textbf{group density}

\[\phi_{r}(S;k) = \frac{|\{(u,v) \in \Pi(S):w_{r}(u,v) \geq k\}|}{|\Pi(S)|}\]

where \(w_{r}\) is the edge weight. At \(k = 1\) this is the share of member pairs connected in \(r\). The level \(k\) matters because a weighted relation can be near saturated as a binary relation while still discriminating at higher weights, and Section 6.5 shows this is the case for co-membership.

\subsection{The Leakage Horizon}

Fix a lead time \(\Delta \geq 0\), in Section 3.1's notation. A forecaster at origin \(t\) predicting events at \(t + \Delta\) sees \(G_{t}\). Define the \textbf{clique contrast} at lead time \(\Delta\),

\[\kappa_{r}(\Delta) = P\left( \phi_{r}(S) = 1 \mid \text{positive} \right) - P\left( \phi_{r}(S) = 1 \mid \text{negative} \right),\]

where a positive is a candidate that accedes at \(t + \Delta\) and a negative is drawn from a declared comparison stratum. The \textbf{leakage horizon} \(\Delta^{\star}\) is the smallest \(\Delta\) at which \(r^{\star}\) is statistically indistinguishable from relations declared not to carry the event.

Three properties of this definition are load bearing and two of them are usually left implicit.

\textbf{The comparison stratum must be declared.} At \(\Delta = 0\), ``no event'' is unambiguous. At \(\Delta > 0\) it is not: a candidate that accedes at \(t + 1\) is a non event for \(\Delta = 3\) and already carries whatever traces the statistic looks for. We compute two strata throughout. The \textbf{never event} stratum contains candidates that never accede to that group anywhere in the observation window, and is the estimand. The \textbf{not at this lead time} stratum is the full complement, which is what a naive implementation produces. The gap between them is contamination introduced by the lead time formulation itself.

\textbf{The reference band must come from relations declared non-leaking on documentary grounds.} No statistic computed from the graph can separate leakage from legitimate signal, because a relation that genuinely informs the label and a relation generated from it both make positives denser than negatives. We therefore fix a set of \textbf{anchor} relations by reading what each source records, before computing anything, and take the band from them. Section 6.6 tests whether that reading is any good.

\textbf{What kind of estimate the rule produces.} It has to be stated, because the obvious name for it promises something else. The specification called the clique-statistic end of the interval a lower bound, and Section 6.1 withdraws the interval. What is left is one number whose direction of error the calibration measures: against a width the generator sets it never falls below the target, and it exceeds it by one or two periods in a third of resamples at the widest setting. That is not the behaviour of a lower bound. A lower bound that exceeds the quantity it bounds has failed, and calling it \emph{conservative} is only defensible once the loss it is conservative under is written down. Treating a contaminated relation as clean corrupts every number computed downstream of it and is not detectable from those numbers. Treating a clean relation as contaminated costs lead times, and the cost is a smaller usable sample. Under that asymmetry an estimator whose errors run upward is the safe one, and under a symmetric loss it would not be. We therefore write \({\widehat{\Delta}}^{\star}\) for what the rule returns, call it a \textbf{conservative horizon estimate}, and keep \(\Delta^{\star}\) for the quantity itself. Where the text discusses the withdrawn interval it says lower and upper \emph{end}, because those name two statistics rather than two bounds.

\subsection{Two Operations on the Focal Group}

Let \(S_{c}\) be the focal group. Two interventions on \(r^{\star}\) are available and they are opposites.

\textbf{Focal exclusion} removes \(E^{r^{\star}} \cap \Pi(S_{c})\), the focal group's own contribution to the relation. It exists so that the group's own edges cannot stand in for the label when estimating a model, and any leakage aware pipeline needs it.

\textbf{Focal isolation} keeps only \(E^{r^{\star}} \cap \Pi(S_{c})\) and discards the rest. It exists because the group's own edges are where administrative back-dating of the group's own event lives: a curator's decision to date a membership from signature rather than operation is a decision about that group's record.

Whether it touches that group's record \emph{alone} is what the operation's freedom from homophily rests on, so it is worth being exact about how far the claim goes. It holds of the focal \emph{event} and it does not transfer to the target \emph{relation}. If the acceding entity and a sitting member are co-members of some other group, the dyad they carry in \(r^{\star}\) belongs to that group's record. If that record is back-dated, the dyad is an administrative trace that no reading of the focal group's register can reach. For such a trace to be correlated with the label the entity's entry into the other group has to fall near its entry into the focal one, which is what accession waves are: expansion rounds, post-independence batches, bloc changes. Section 6.3 makes the wave a controlled parameter and measures what it costs. Under a single dating convention the answer is nothing, because the other group's record is back-dated by the same width and puts nothing into the relation that the focal record does not already carry. Under mixed conventions, which is corpus A exactly, the answer is up to nine periods of reach that a focal reading does not see, and the remedy is to read every register that feeds \(r^{\star}\) rather than the focal group's alone.

The two operations serve opposite purposes and Section 6.4 shows that a suite offering only exclusion cannot measure the quantity it exists to measure whenever \(r^{\star}\) carries contributions from anything other than the focal group. In a curated corpus it always does.

\subsection{What Is Measured, and What Is Entailed}

Several of the results below are zeros, and a zero can arise in two ways that mean opposite things. A statistic can return zero because the quantity it looks for is absent, which is a measurement. It can also return zero because the register cannot express the quantity at all, in which case the zero follows from the construction and carries no information about the corpus. The distinction is easy to state and easy to lose, and losing it is what turns a bounded claim into an inflated one, so we fix it here once and refer back rather than restating it at each occurrence. In the vocabulary of measurement modelling \citep{jacobs} this is the gap between a construct and its operationalisation, and each row of Table 1 records which side of that gap a number sits on.

\textbf{What a pooled contrast measures.} A contrast taken over the whole target relation is not a trace of the event, whether it is the clique contrast of Section 3.2 or the mean density contrast of Table 7. It is a sum of contributions, of which the trace is one:

\[\kappa_{r^{\star}}(\Delta)\mspace{6mu} = \mspace{6mu}\underset{\begin{matrix}
\text{trace on the} \\
\text{focal record}
\end{matrix}}{\underbrace{\tau(\Delta)}}\mspace{6mu} + \mspace{6mu}\underset{\begin{matrix}
\text{trace on} \\
\text{other records}
\end{matrix}}{\underbrace{\lambda(\Delta)}}\mspace{6mu} + \mspace{6mu}\underset{\text{homophily}}{\underbrace{\eta(\Delta)}}\mspace{6mu} + \mspace{6mu}\underset{\text{recurrence}}{\underbrace{\rho(\Delta)}}\mspace{6mu} + \mspace{6mu}\varepsilon.\]

This is a decomposition of sources rather than an algebraic identity, and it is useful because each term has a section:

\begin{itemize}
\tightlist
\item
  \(\tau\), the administrative trace on the focal group's own record, is the estimand focal isolation reads, and it is the only term it reads.
\item
  \(\lambda\), the administrative trace of the same event carried by \emph{another} group's record, is separated from \(\eta\) because Section 6.3 measures it and because it is administrative rather than a similarity. It is zero when every register in the corpus dates by one convention and it is not zero when they differ, and no reading of the focal record recovers it. It is mechanism (k).
\item
  \(\eta\), persistent similarity between the candidate and the members, is mechanism (h). It is present at every lead time, including lead times far beyond any dating convention, and Section 6.2 measures how large the resulting error is as a function of homophily strength on a generator where \(\tau\) is set by hand.
\item
  \(\rho\), prior membership spells, is mechanism (i). It \emph{grows} with lead time and is removed by conditioning on first entry.
\item
  \(\varepsilon\) is sampling noise, and Section 6.1's negative control is where its false-positive rate is measured.
\end{itemize}

No contrast of this form separates \(\tau\) from \(\eta\), because a candidate that is denser to the members because the register back-dated its accession and a candidate that is denser because it was always similar to them enter \(\kappa\) identically. Section 6.2 turns that from a statement into a measured bias.

\textbf{What focal isolation measures, and under what assumption.} Focal isolation asks whether the acceding entity is on the focal group's own record \(\Delta\) periods before the coded event date. It identifies the back-dating width \textbf{if and only if the register dates the relation and the event from separate fields.} Where it does, as in Section 5.3's trade agreement register, the presence curve is the width and the measurement is exact. Where the register codes one date per membership and derives the event from a transition in that same field, the curve is 1 at \(\Delta = 0\) and 0 beyond \textbf{by construction}: there is no second date for the two to disagree about, so a zero residual is entailed and not measured. Corpus A is of the second kind for 94 percent of its events, and Section 8 says what to do about it.

\begingroup\small
\begin{longtable}{@{}
  >{\raggedright\arraybackslash}p{(\linewidth - 4\tabcolsep) * \real{0.4620}}
  >{\centering\arraybackslash}p{(\linewidth - 4\tabcolsep) * \real{0.1485}}
  >{\raggedright\arraybackslash}p{(\linewidth - 4\tabcolsep) * \real{0.3895}}@{}}
\caption{Every headline quantity, and whether it is a measurement, an entailment of the construction, or a validation of the pipeline. Entailments are not weaker results; they are results about the register rather than about the corpus, and reading one as the other is the error this table exists to prevent.}\label{tab:1}\\
\toprule
\begin{minipage}[b]{\linewidth}\raggedright
\textbf{Quantity}
\end{minipage} & \begin{minipage}[b]{\linewidth}\centering
\textbf{Where}
\end{minipage} & \begin{minipage}[b]{\linewidth}\raggedright
\textbf{Status}
\end{minipage} \\
\midrule
\endfirsthead
\toprule
\begin{minipage}[b]{\linewidth}\raggedright
\textbf{Quantity}
\end{minipage} & \begin{minipage}[b]{\linewidth}\centering
\textbf{Where}
\end{minipage} & \begin{minipage}[b]{\linewidth}\raggedright
\textbf{Status}
\end{minipage} \\
\midrule
\endhead
\bottomrule
\endlastfoot
The estimate recovers \(\beta + 1\), 6 of 6 settings, 67 to 100 percent of replicates & 6.1, Table 3 & \textbf{measurement} against a width set by the generator \\
Synthetic false-positive rate at \(h = 0\), \(\beta = 0\) & 6.1 & \textbf{measurement}, the negative control \\
Identification slope of each estimator on \(\beta\) & 6.2, Table 4 & \textbf{measurement} against a width set by the generator \\
Censoring rate behind each slope, and the same slope over a range twice as wide & 6.2, Tables 4 and 5 & \textbf{measurement} about the estimator, not about a corpus \\
Off-focal administrative reach under one dating convention & 6.3, Table 6 & \textbf{measurement}, and it is zero \\
Off-focal administrative reach under mixed conventions & 6.3, Table 6 & \textbf{measurement} against a width set by the generator \\
PTA back-dated increment, 6.8 percent of the signature-dated relation & 5.3 & \textbf{measurement}, both dates coded \\
PTA focal record, presence 1 inside the coded gap and 0 outside & 6.4 & \textbf{entailment}, and reported as a pipeline validation \\
PTA pooled contrast 0.1023 & 6.4 & \textbf{measurement} \\
Corpus A exactly dated stratum, 1.000 falling to 0.000 & 6.5 & \textbf{entailment}, one coded date per membership \\
No stratum of corpus A supports both instruments as measurements & 6.5 & \textbf{measurement} about the corpus, and the reason the falsification rests elsewhere \\
Corpus A declared-resolution stratum, 0.1419 at one period & 6.5 & \textbf{measurement}, the register's own block resolution \\
Corpus A focal contribution, 6.5 percent of the relation at \(k = 25\) & 6.5 & \textbf{measurement} \\
Corpus A pooled contrast, 0.3114 falling to 0.1490 & 6.5, Table 7 & \textbf{measurement} \\
Corpus A four-period decay attributable to recurrence & 6.5, Figure 5 & \textbf{measurement} \\
Clique contrast 0.0000 above eighty members & 6.6, Table 8 & \textbf{measurement} with a combinatorial ceiling \\
Corpus B channel ranking, first, second and third & 6.6 & \textbf{measurement}; whether it is also a prospective test is not checkable, see 6.6 \\
The interval's upper end unreachable, or below its lower end & 6.1 & \textbf{measurement} about the definition, not about a corpus \\
\end{longtable}
\endgroup

\section{The Diagnostic Suite}

The suite has four components, and the fourth is the one Section 6.4 shows the design cannot do without. Figure 1 gives the pipeline.

\textbf{U1, the horizon curve.} For each lead time \(\Delta\) in a declared range, compute the clique contrast for \(r^{\star}\) and for every anchor, on both comparison strata, at every level \(k\) of the declared ladder. The estimate \({\widehat{\Delta}}^{\star}\) is the smallest \(\Delta\) at which the clique contrast for \(r^{\star}\) falls inside the anchor band. Section 3.2 states what the number is and is not: a conservative estimate under an asymmetric loss, and not a lower bound on \(\Delta^{\star}\).

A second statistic was specified alongside it, to close the interval from above, and we report it as non-identified rather than as an estimate. As specified, the upper end was the smallest \(\Delta\) at which \(r^{\star}\) both loses top rank on clique contrast and has its single-feature discrimination fall inside the band. The pairing was meant to be forced by a property of the two statistics. The clique contrast raises no false flags but only fires above a high threshold, while rank position and single-feature discrimination fire earlier and are not specific, so one should declare the horizon closed early and the other late. Two measurements say that pairing does not hold. On synthetic data the rank condition blocks the upper end in half the settings, because a structurally similar anchor keeps contesting the top position. Dropping the rank condition makes it fire everywhere but land on \({\widehat{\Delta}}^{\star}\) in every single setting, so the interval it forms has width zero. On corpus A at the declared primary level the two statistics disagree in the wrong direction at every lead time examined. The clique contrast sits outside its band throughout, so the specific statistic calls the relation distinguishable everywhere. The single-feature discrimination sits inside its band throughout, so the sensitive statistic calls it indistinguishable everywhere. Where both exist the second falls below the first. \textbf{The interval is therefore withdrawn and the suite reports one number.} We keep the definition in print because a rule that failed is evidence about the design, and because a prediction bounded above cannot be tested against a bound that does not exist.

\begin{quote}
\textbf{The level ladder, stated as a convention.} Three rules, and the third is the one an independent re-implementation could not recover from the first specification.
\end{quote}

\begin{enumerate}
\def\labelenumi{\arabic{enumi}.}
\tightlist
\item
  \textbf{The ladder applies to} \(r^{\star}\) \textbf{only.} It exists because a weighted relation can be saturated as a binary relation and still discriminate at higher weights, and co-membership's weight is a count of shared organisations. \(r^{\star}\) is read at level 1 and at the declared primary level.
\item
  \textbf{The anchors are read as presence relations.} Their weights are not counts: trade flows carry log volumes from \(- 3.98\) to \(3.54\) and voting agreement lies strictly inside the unit interval.
\item
  \textbf{``Level 1'' means presence, not weight at least one.} A literal threshold at weight one would empty the voting relation and discard half of trade flows.
\end{enumerate}

\begin{quote}
The third rule is the one an independent re-implementation could not recover from the prose: written from the specification, it read level 1 literally and disagreed with the original by exactly 1.0 on those two relations. Section 10 reports that cross-check.
\end{quote}

\textbf{U2, the lead-time-resolved shortcut probe.} Fit a probe using only relational densities and clique indicators, no learned representation, at each lead time, cross-fitted with folds grouped on the resampling unit. Report the discrimination curve and its decay. Report a linear probe alongside a non-linear one, for a reason Section 7 gives.

\textbf{U3, the exclusion contrast.} Refit with \(r^{\star}\) excluded and report the change. This measures how much of a fitted model's discrimination depends on the target relation, as distinct from how much of the target relation's density is event related.

\textbf{Focal isolation.} For each event, ask whether the acceding entity is already on the focal group's own record \(\Delta\) periods before the coded event date. Under honest dating this is 1 at \(\Delta = 0\) and 0 beyond, and any excess is back-dating of exactly that width. This is the component the design cannot do without, and Section 6.4 quantifies what its absence costs. Applied to a corpus it is run over every group's own record, not only over the group whose event is being audited, and Section 6.3 measures why the difference matters: where dating conventions differ between groups, the narrower reading undershoots and the wider one does not.

\begin{figure}[!htbp]
\centering
\includegraphics[width=\textwidth]{figures/figure-1.pdf}
\caption{\textbf{The diagnostic suite, and the two questions it separates.} Schematic, no data.}
\end{figure}

\section{Experimental Design}

We use four testbeds. Two are ground-truth settings where the answer is known independently of the diagnostic, and two are field corpora where it is not.

\subsection{Declaration Protocol}

Every choice that could be made after seeing a result was fixed in writing beforehand and is reported whether or not it flattered the outcome. Where a declared choice was later changed, the body carries a tag of the form and \textbf{Appendix B} says what the change was and what forced it, so that an amendment is visible without a reader having to sort a conclusion from its history at every occurrence. This includes the target relation and the anchor set for each corpus, the level ladder for weighted relations, the comparison strata, the resampling unit, the covariate set for the survival models, and the lead times. Where a pre-specified rule returned a result we then judged uninterpretable, we report the rule's output and the reason for rejecting it, rather than substituting a different rule.

The protocol is meant to carry the weight of Section 6.6's reading of ten channels, which placed the three contaminated ones first, second and third, so ``written beforehand'' has to be checkable rather than stated. It is checkable for one of the two corpora. The declaration is therefore released as a separate citable artefact, cited under data and code availability, with the text fixed at the version the reported statistics were computed against and a checksum for each file.

\textbf{One half of the declaration is timestamped and the other is not.} The weaker half is the one the prospective test rests on. Corpus A's declaration is a pre-registration in the strict sense: it was committed to version control before any data was materialised, amended twice within a day, and frozen, and the three commits carry their dates. Everything the protocol names for corpus A is in it, including the two predictions it went on to falsify. Corpus B's channel declaration is not. The ten STRING channels were assigned to target, contaminated and anchor in the docstring of the analysis script, which is where we wrote them and is our account of the order. That file was not committed before it was run, and the pre-registration does not name the channels at all. So for corpus B the repository cannot corroborate that the reading preceded the statistic.

Depositing the artefact now does not repair this, and saying that it does would be the error Section 9 is about: a timestamp obtained after a statistic is not evidence about the order in which the two were produced. What the deposit does is fix the text so it cannot drift, and record the distinction in the artefact itself, and the two files it contains are marked timestamped and not timestamped accordingly. Section 6.6 states the same limitation where the test is reported, because a reader who only reads that section should not have to come back here for it.

Three of the changes reported in this paper are revisions \emph{to} the protocol rather than results produced under it, and Appendix B marks which three. All three are marked as such where they appear, because a pre-registration that is silently amended is worse than none. The deposited artefact carries the pre-registration as frozen, so what was changed is checkable and not only reported.

\subsection{Synthetic Testbed}

A generator produces a group corpus with two controlled parameters: \textbf{completeness} \(\alpha\), the share of the target relation's edges that are present, and \textbf{back-dating} \(\beta\), the number of periods before the event at which the focal group's edges appear. The corpus has 200 nodes and 90 accession events per setting, and completeness tracks \(\alpha\) to within 0.0004.

\textbf{What a calibration testbed has to get right.} The density of \(r^{\star}\) must not move with the parameter being calibrated, or the estimator is graded against a relation that got denser rather than against a width that got wider, and every horizon curve slopes for a reason unrelated to leakage. Here the density is \textbf{exactly} invariant, one value across every cell of the grid, because the edge budget is a fixed integer count of dyads rather than a sampled quantity, and the invariance is checked before any estimator is graded rather than reported afterwards. The tempting alternative is worse than doing nothing: \textbf{letting the budget grow with} \(\beta\) \textbf{restores invariance on paper and destroys it in fact.} Back-dating enlarges the pool of group-induced dyads. A budget that grows to cover them keeps the nominal density fixed while the group-induced share of it rises with the parameter, which is precisely the confound the constraint exists to remove. The budget is therefore sized once, from the measured worst case over the whole sweep, and the peak eligible-dyad count that justifies it is an artefact rather than a comment: 1,794 dyads of 19,900 pairs at the widest setting used. Section 6.3's sweep sizes its own budget the same way and reports the headroom it had.

\textbf{A second generator adds the axes the first holds fixed.} The generator above cannot exhibit the paper's deepest claim, because its similarity structure is fixed. The background of every relation is drawn by latent proximity and the acceding entities are always the group's nearest non-members, so there is no axis along which homophily can be shown to grow. Section 6.2 uses an extension with four further controls, and the first generator is left untouched so that Table 3's calibration cannot move.

\begin{itemize}
\tightlist
\item
  \textbf{Homophily} \(h\), an exponent on latent similarity in every sampling weight: in the edge background, in the founding membership, and in the pool the acceding entities are drawn from. At \(h = 0\) the background is a uniform random sample and the acceding entity is a uniform random non-member, so positives and negatives differ only through the event. At \(h = 1\) similarity enters linearly and above that it concentrates.
\item
  \textbf{A width distribution.} A sharp \(\beta\) is a regime an estimator can hardly fail in, so the width is redrawn per event. We draw it from the PTA register's own coded gap distribution of Section 5.3, which makes the widths ones a real curator produced rather than ones we invented.
\item
  \textbf{Dating noise}, a share of events whose \emph{coded} date is displaced from the true one. The register is wrong rather than early, which is the synthetic counterpart of corpus A's declared-resolution stratum.
\item
  \textbf{Whether the anchors share the latent space.} Drawn from the same latent trajectory as the target relation they carry the same homophily, which is corpus A; drawn from an independent one they carry none of it, which is corpus B. One switch produces both, so the difference between the two field corpora becomes a setting rather than a conjecture.
\item
  \textbf{A curatorial wave, and the convention the other register dates by.} A share of accession events arrive with the acceding entity also entered into a second group that already overlaps the focal one, at the same true period. The second register is back-dated by its own law, which may be the focal register's or a different one. This is the axis Section 6.3 needs and the generator had no way to express. With one convention the wave puts nothing into \(r^{\star}\) that the focal record does not carry; with two it puts in a trace no reading of the focal record can reach. The generator records, per event, which record carries the trace and how far it reaches, so the two are separable rather than inferred.
\end{itemize}

The generator also records, for every accession event, the reach of the administrative trace and the register that carries it, split three ways: the focal group's own record, a second record the event's own wave back-dated, and the ambient content that any group's back-dating puts into \(r^{\star}\) whether or not the focal event happened. The third is reported against matched negatives, because an administrative dyad that positives and negatives carry alike is not a trace of the event and reading it as one would be the error this paper is about.

Two registers are emitted rather than one, because the two instruments read different objects. The honest register records an entity as a member from the period it truly joined, and U1 has to be built on that one: the risk set is states not already members, so a back-dated register would drop exactly the leaking rows as ineligible. The recorded register is what a curator wrote, with every event appearing \(\beta\) periods early, and focal isolation reads that one. This is Section 5.3's situation exactly, a relation dated from signature against an event dated from entry into force, which is what lets the two instruments be compared on one corpus against one known truth. The relation's density is again exactly invariant across every cell, at a level sized from the measured worst case rather than guessed: the widest setting used, the PTA gap distribution with its tail at nine periods, peaks at 1,794 eligible dyads of 19,900 pairs.

\subsection{A Field Corpus with Coded Ground Truth}

A register of preferential trade agreements \citep{desta} codes, for each agreement and each acceding party, both a signature year \(s\) and an entry into force year \(f\). An observation graph built from signature dated membership therefore carries a dyad exactly \(f - s\) periods before the legal event. The width is printed in the source. Over the observation window there are 759 accession entries with a coded force year, with gap median 1, mean 1.71, and 39.4 percent at two or more periods. The signature dated relation has 204,747 edge-years against the force dated relation's 190,781, so the two differ by 13,966 edge-years, which is 6.8 percent of the signature dated relation. The convention adds dyad-years rather than moving them: 13,967 dyad-years appear only in the signature dated relation and exactly one appears only in the force dated one, so the difference is a back-dated increment and not a shift. \textbf{This 6.8 percent is a property of this register and is not carried to corpus A}, whose own focal contribution is computed in Section 6.5 and is a different number at every level.

\subsection{Field Corpus A: International Coalitions}

The primary corpus comprises 562 entities, 490 intergovernmental organisations and 72 multilateral treaties, over 1946 to 2014, carrying 10,328 events: 8,062 accessions, 1,589 exits, 492 foundings and 185 terminations. Membership occupies 505,400 entity-year rows over 202 states. Events fall in three dating strata: 9,581 exact, 646 declared resolution, where the source states its own dates are accurate only within five years, and 101 rule assigned, where the coding rule computes a date from up to a decade of subsequent data. Terminations decompose into 101 dissolutions, 32 replacements, 22 absorptions and 30 whose cause the source does not state.

The observation graph carries five relations: co-membership, alliance ties from a treaty obligations register \citep{atop}, trade flows, voting agreement in the United Nations General Assembly derived from dynamic ideal points \citep{bailey}, and militarised disputes \citep{mid5}. Analysis is confined to 1965 to 2014, the maximal window over which the set of non-empty relations does not change, because a coverage break inside a pooled metric produces gains that appear in no sub-window.

Resampling uses \textbf{membership equivalence families} rather than entities. Two entities whose member sets coincide for a sustained period are not independent observations, and the corpus contains 464 families of which 43 hold more than one entity.

\subsection{Field Corpus B: Protein Complexes}

The replication corpus comprises 2,540 curated protein complexes against 2,540 random and 5,080 hard negatives over 19,486 proteins, with a ten-channel interaction network as the relations. It was chosen for one property: its group size distribution is the inverse of corpus A. Median complex size is 3, mean 5.09, maximum 129, with 1.57 percent of positives above twenty members and 0.16 percent above eighty. Corpus A's positive strata average 73.5 members, median 60, with 41.1 percent above eighty. Negatives are already size-matched to positives, median 3 and mean 5.09 on both sides.

The target relation is the channel that imports the curated complex registers, which is the label source itself, by the same logic that makes co-membership the target relation in corpus A. Two further channels were declared \textbf{contaminated and therefore ineligible as anchors} before any statistic was computed: one aggregates text-mined co-occurrence from papers that describe complexes, and one aggregates experimental evidence including co-immunoprecipitation, which detects complexes. The five remaining channels were declared anchors: messenger RNA co-expression, two transferred variants, gene fusion events, phylogenetic co-occurrence, and sequence homology. None of these can record a curator's decision that a set of proteins forms a complex.

\textbf{Negative construction, stated in full} \citep{poursafaei}, because the level of every contrast we report depends on it. Each positive complex generates three negatives: two \emph{hard} and one \emph{random}. A random negative is \(k\) proteins drawn uniformly without replacement from the universe of 19,486 proteins carrying at least one channel edge, where \(k\) is the positive's size, so it is size-matched by construction. A hard negative keeps the positive's members with \(b = min(2,k - 2)\) of them dropped and \(b\) proteins substituted from the same universe. Complexes below three mapped subunits are excluded, so the smallest groups carry three members and three dyads. \textbf{No existing edge is excluded and no collision with a real complex is checked.} That matters because the alternative would make part of the reported contrast structural rather than measured. It is verifiable rather than asserted: if edges had been avoided the negatives' mean density in the target relation would be zero, and it is 0.00132 against the relation's own density of 0.00115. Table 8 reports the resulting clique counts as integers rather than as rates.

Table 2 states the two field corpora side by side, and the flow beneath it states every filter that separates one reported count from another, so that a reader can check a difference rather than reconstruct it.

\begin{table}[!htbp]
\centering
\small
\begin{tabular}{@{}
  >{\raggedright\arraybackslash}p{(\linewidth - 4\tabcolsep) * \real{0.3645}}
  >{\centering\arraybackslash}p{(\linewidth - 4\tabcolsep) * \real{0.3468}}
  >{\centering\arraybackslash}p{(\linewidth - 4\tabcolsep) * \real{0.2887}}@{}}
\toprule
\begin{minipage}[b]{\linewidth}\raggedright
\end{minipage} & \begin{minipage}[b]{\linewidth}\centering
\textbf{Corpus A, coalitions}
\end{minipage} & \begin{minipage}[b]{\linewidth}\centering
\textbf{Corpus B, protein complexes}
\end{minipage} \\
\midrule
Groups & 562 entities & 2,540 complexes \\
Nodes & 202 states & 19,486 proteins \\
Positive events & 8,062 accessions & 2,540 complexes \\
Negatives & declared strata, subsampled & 2,540 random, 5,080 hard \\
Relations & 5 & 10 \\
Time axis & 1946 to 2014, annual & none \\
Median group size, all group periods & 5 alliances, 16 organisations & 3 \\
Median size of the group an event enters & 60 & 3 \\
Share above 80 members & 41.1 percent of events & 0.16 percent of complexes \\
Target relation density & 0.9954 & 0.001153 \\
Highest relation density & 0.9954 & 0.001887 \\
\bottomrule
\end{tabular}
\caption{Field corpora. The size distributions are inverses of each other, which is why corpus B is the replication.}\label{tab:2}
\end{table}

\textbf{Corpus A, from the register to each reported count.} Every arrow is a filter and every filter is named.

{\small
\begin{verbatim}
562 entities
  -> [observation interval overlaps 1965 to 2014] 528 groups, the population left
     truncation threatens, which is the count Section 9 reports
  -> [an active membership row inside the window] 527 groups, the panel the
     termination hazard is estimated on, which is the count Section 6.7 and Appendix
     A report
185 coded terminations
  -> [inside the window] 151
  -> [on a group in the estimated panel] 150 events, Appendix A
8,062 accession events
  -> [inside the window] 6,597 events, Table 8 and Figure 7
  -> [one accession per dyad, since a dyad leaves the risk set when it accedes]
     6,501 event dyads, Table 9; 96 repeat accessions of 94 dyads are collapsed onto
     the first, 1.5 percent
7,615 exactly dated accessions, whole register
  -> [first entry, which removes prior membership spells] 7,146 events at lead time
     zero, Section 6.5
  -> [the focal group has a member set in the preceding period] 7,084 events at lead
     time one, which is what the interval on that stratum is taken over; the 62
     dropped are accessions in a group's founding period, where there is no earlier
     record to be absent from
447 declared-resolution accessions
  -> [first entry] 310 events, and none is lost at lead time one
\end{verbatim}
}

The whole chain is recomputed by a script that reaches each count through the same code path the reporting script uses, so this is a second reading of the numbers rather than a restatement of them, which is the fifth question of Section 9 applied to a count instead of to a density.

\section{Results}

The results run from the testbed, where the width is a parameter we set and the answer is therefore known, to two field corpora, where the width is a property of a register we did not build, and end at what the horizon costs a model trained across it.

\subsection{Calibration Against a Width the Generator Sets}

Table 3 reports the pre-specified rule applied to the synthetic testbed at six settings of the back-dating parameter. The estimand implied by the definition is \(\beta + 1\): traces are present for \(\Delta \leq \beta\), so the smallest trace free lead time is \(\beta + 1\).

\begin{table}[!htbp]
\centering
\small
\begin{tabular}{@{}
  >{\raggedright\arraybackslash}p{(\linewidth - 10\tabcolsep) * \real{0.2898}}
  >{\centering\arraybackslash}p{(\linewidth - 10\tabcolsep) * \real{0.1273}}
  >{\centering\arraybackslash}p{(\linewidth - 10\tabcolsep) * \real{0.0808}}
  >{\centering\arraybackslash}p{(\linewidth - 10\tabcolsep) * \real{0.1399}}
  >{\centering\arraybackslash}p{(\linewidth - 10\tabcolsep) * \real{0.1939}}
  >{\centering\arraybackslash}p{(\linewidth - 10\tabcolsep) * \real{0.1682}}@{}}
\toprule
\begin{minipage}[b]{\linewidth}\raggedright
\textbf{Width regime}
\end{minipage} & \begin{minipage}[b]{\linewidth}\centering
\textbf{Target}
\end{minipage} & \begin{minipage}[b]{\linewidth}\centering
\[{\widehat{\mathbf{\Delta}}}^{\mathbf{\star}}\]
\end{minipage} & \begin{minipage}[b]{\linewidth}\centering
\textbf{On target}
\end{minipage} & \begin{minipage}[b]{\linewidth}\centering
\textbf{Replicates on target}
\end{minipage} & \begin{minipage}[b]{\linewidth}\centering
\textbf{95 percent interval}
\end{minipage} \\
\midrule
sharp, \(\beta = 0\) & 1 & 1 & yes & 100 percent & 1 to 1 \\
sharp, \(\beta = 1\) & 2 & 2 & yes & 100 percent & 2 to 2 \\
sharp, \(\beta = 2\) & 3 & 3 & yes & 91 percent & 3 to 5 \\
sharp, \(\beta = 3\) & 4 & 4 & yes & 92 percent & 4 to 6 \\
sharp, \(\beta = 4\) & 5 & 5 & yes & 67 percent & 5 to 7 \\
sharp, \(\beta = 5\) & 6 & 6 & yes & 92 percent & 6 to 8 \\
PTA gaps, graded against the realised maximum of 4 & 5 & 5 & yes & not resampled & \\
PTA gaps, graded against the realised 90th percentile of 4 & 5 & 5 & yes & not resampled & \\
PTA gaps, graded against the realised median of 1 & 2 & 5 & no, high by three & not resampled & \\
\bottomrule
\end{tabular}
\caption{Calibration of the conservative estimate against a known back-dating width, primary comparison stratum, 90 accession events per setting. Columns five and six resample coalitions within each setting, 1,000 replicates.}\label{tab:3}
\end{table}

The heterogeneous rows are one corpus per homophily setting rather than a sweep, and the estimate is 5 in all twelve of them, so they are reported as a single block. The rule recovers the maximum width and not the median one, which is a definitional result rather than a lucky one: the clique contrast leaves the anchor band only when \textbf{no} event still carries a trace, and the last event to stop carrying one is the widest.

The estimate lands on the target in \textbf{6 of 6} settings, and Figure 2 shows the six estimates against the target line. The resampling column is what a point estimate hides. At four of the six settings the rule is essentially deterministic at this sample size. At \(\beta = 4\) it returns the target in two replicates out of three and overshoots by one or two in the rest, and it never undershoots in any replicate at any setting. Both halves of that sentence matter. \textbf{Read as a lower bound the rule fails 33 percent of the time at} \(\beta = 4\)\textbf{.} A lower bound that comes back above the quantity it bounds is not conservative, it is wrong. The synthetic testbed is the one place where that can be checked, because \(\Delta^{\star} = \beta + 1\) is known exactly there. Read as the conservative estimate Section 3.2 defines, the same number is a 67 percent coverage rate on the truth with every miss on the side that discards data rather than admits it. We report the second reading, and the coverage rate beside the six ticks, because the ticks are point estimates and the rate is what a reader choosing how many lead times to discard would actually want.

\textbf{The negative control.} At homophily \(h = 0\) and \(\beta = 0\) there is nothing to find, so an instrument that reports a horizon there is reporting noise. All four estimators of Section 6.2, and the pre-specified rule among them, return the correct value of 1: no false positive. This cell is the synthetic counterpart of the one relation in corpus A that the diagnostic separates from the target: militarised disputes, flat at 0.004 at every lead time in Table 7. The two are the same statement about the same instrument, made once where the truth is set and once where it is not.

\textbf{The interval's upper end is not identified.} The repair does not repair it. As specified it fired in \textbf{3 of 6} settings, at \(\beta = 0\), \(1\) and \(2\), and was unreachable in the rest, because it requires \(r^{\star}\) to lose top rank and a structurally similar anchor keeps contesting the top position. Over the wider grid of Section 6.2 it fired in 20 of 72 cells. Dropping the rank condition, so that it is a band crossing alone, makes it fire in 72 of 72 cells and land on \({\widehat{\Delta}}^{\star}\) \textbf{in every one of them}, so the interval it forms has width zero and the second statistic contributes nothing. On corpus A at the declared primary level the repaired version fires at one period while \({\widehat{\Delta}}^{\star}\) never fires at all inside ten, and at the one level where both exist the upper falls below the lower. The two statistics do not order the way the bracket requires: the clique contrast sits outside its band at every lead time from one to ten, and the single-feature discrimination sits inside its band at every one of them. The ordering has failed twice on this corpus and the second failure is what identifies the cause. A superseded build produced an inverted interval on the alliance subcorpus, a lower end of two above an upper end of one, at a point when the affected stratum had lost 47 percent of its positives, which looks like a power problem. The same inversion appears on the current corpus in the organisation stratum, at full sample, at a level where both ends exist, and the synthetic grid shows the repaired version sitting on \({\widehat{\Delta}}^{\star}\) in all 72 cells. A failure that survives a corpus rebuild, a change of stratum and a change of sample size is a property of the two statistics rather than of a sample. We therefore report \textbf{instrumental non-identification of the upper end}, as a result about the definition rather than about any corpus, and we withdraw the interval in favour of one conservative estimate. A prediction bounded above cannot be falsified by a bound that does not exist, and it is not rescued by one that is \({\widehat{\Delta}}^{\star}\) under another name.

\begin{figure}[!htbp]
\centering
\includegraphics[width=\textwidth]{figures/figure-2.pdf}
\caption{\textbf{How the estimate is formed, and that it lands on the target six times out of six.} Synthetic testbed, six settings of the back-dating width, 200 nodes and 90 accession events per setting, clique contrast at level \(k = 1\) on the never-event comparison stratum, lead times 1 to 8.}
\end{figure}

\subsection{Homophily Against Back-dating}

Section 3.4 states that a pooled contrast sums an administrative trace and a persistent similarity, and that no contrast of that form separates them. On a field corpus that claim cannot be tested, because the trace is the unknown. On the generator of Section 5.2 it can: the width is set, homophily is a dial, and an estimator's error is therefore a function of the dial. We sweep six homophily settings against six widths in each of the two anchor regimes, 72 corpora, and grade four estimators against the same target \(\beta + 1\).

Two summaries carry the result. The \textbf{mean error} is how far an estimator sits from the truth, averaged over the six widths at one homophily setting. The \textbf{slope on} \(\beta\) is the least-squares slope of the estimate on the true width across those six settings: a slope of 1 means the estimator tracks the truth, and a slope of 0 means its output is statistically unrelated to the quantity it is estimating. An estimate may be biased and still informative; an estimate with slope 0 is not.

\begin{table}[!htbp]
\centering
\small
\begin{tabular}{@{}
  >{\raggedright\arraybackslash}p{(\linewidth - 10\tabcolsep) * \real{0.1448}}
  >{\centering\arraybackslash}p{(\linewidth - 10\tabcolsep) * \real{0.0620}}
  >{\centering\arraybackslash}p{(\linewidth - 10\tabcolsep) * \real{0.1876}}
  >{\centering\arraybackslash}p{(\linewidth - 10\tabcolsep) * \real{0.1876}}
  >{\centering\arraybackslash}p{(\linewidth - 10\tabcolsep) * \real{0.2091}}
  >{\centering\arraybackslash}p{(\linewidth - 10\tabcolsep) * \real{0.2090}}@{}}
\toprule
\begin{minipage}[b]{\linewidth}\raggedright
\textbf{Anchors}
\end{minipage} & \begin{minipage}[b]{\linewidth}\centering
\[\mathbf{h}\]
\end{minipage} & \begin{minipage}[b]{\linewidth}\centering
\textbf{Focal isolation}
\end{minipage} & \begin{minipage}[b]{\linewidth}\centering
\textbf{Clique estimate}
\end{minipage} & \begin{minipage}[b]{\linewidth}\centering
\textbf{Residual rule}
\end{minipage} & \begin{minipage}[b]{\linewidth}\centering
\textbf{Unbanded rule}
\end{minipage} \\
\midrule
share the latent space & 0.0 & 0.00 / 1.00 / 0 & 0.00 / 1.00 / 0 & 0.00 / 1.00 / 0 & 0.00 / 1.00 / 0 \\
& 0.5 & 0.00 / 1.00 / 0 & 0.00 / 1.00 / 0 & 0.00 / 1.00 / 0 & 0.00 / 1.00 / 0 \\
& 1.0 & 0.00 / 1.00 / 0 & 0.00 / 1.00 / 0 & 0.00 / 1.00 / 0 & 0.00 / 1.00 / 0 \\
& 2.0 & 0.00 / 1.00 / 0 & 0.00 / 1.00 / 0 & 0.00 / 1.00 / 0 & +4.33 / -0.03 / 4 \\
& 3.0 & 0.00 / 1.00 / 0 & 0.00 / 1.00 / 0 & 0.00 / 1.00 / 0 & +5.50 / 0.00 / 6 \\
& 4.0 & 0.00 / 1.00 / 0 & 0.00 / 1.00 / 0 & 0.00 / 1.00 / 0 & +5.50 / 0.00 / 6 \\
independent & 0.0 & 0.00 / 1.00 / 0 & 0.00 / 1.00 / 0 & 0.00 / 1.00 / 0 & 0.00 / 1.00 / 0 \\
& 0.5 & 0.00 / 1.00 / 0 & 0.00 / 1.00 / 0 & 0.00 / 1.00 / 0 & 0.00 / 1.00 / 0 \\
& 1.0 & 0.00 / 1.00 / 0 & 0.00 / 1.00 / 0 & 0.00 / 1.00 / 0 & 0.00 / 1.00 / 0 \\
& 2.0 & 0.00 / 1.00 / 0 & 0.00 / 1.00 / 0 & +2.00 / +0.54 / 1 & +4.33 / -0.03 / 4 \\
& 3.0 & 0.00 / 1.00 / 0 & 0.00 / 1.00 / 0 & +5.33 / -0.03 / 6 & +5.50 / 0.00 / 6 \\
& 4.0 & 0.00 / 1.00 / 0 & 0.00 / 1.00 / 0 & +5.50 / 0.00 / 6 & +5.50 / 0.00 / 6 \\
\bottomrule
\end{tabular}
\caption{Identification. Mean error, slope on the true width, and censoring, per estimator, per homophily setting, in each anchor regime.}\label{tab:4}
\end{table}

\textbf{Homophily alone destroys identification in the pooled statistic.} Read the last column down. Up to \(h = 1\) the unbanded rule is exact. At \(h = 2\) its mean error is four and a third periods and its slope has collapsed to \(- 0.03\); at \(h = 3\) and above it reports the largest lead time examined at every one of the six widths, so within this range its output is a constant. Nothing about the corpus changed except how strongly entities that are close in the latent space are tied to one another. The instrument reads that closeness as back-dating, and the exchange rate is steep: between \(h = 1\) and \(h = 3\) the reported horizon moves from the truth to the end of the range while the truth stays put.

\textbf{Widening the range.} A constant at the top of a range is a fact about the range as well as about the estimator. The censoring column says exactly where this bites. At \(h \geq 3\) the collapsed cells are censored in all six of their widths, so their slope of zero is arithmetically forced and, taken alone, says nothing. At \(h = 2\) the collapse is partial and the column is worth reading rather than skipping. The unbanded rule is censored in four widths of six and its slope over the two that are not is \(+ 0.50\). The residual rule with independent anchors is censored in one of six, with an uncensored slope of \(+ 0.64\) against the \(+ 0.54\) the cell reports. Those two cells are transitional in fact and not only in appearance. Table 5 reruns those cells over lead times to twenty, two and a half times the original range, and the result is worse for the estimator rather than better. The estimates move with the range, from a constant 9 to values between 11 and 21, which confirms the censoring. The dependence on the truth does not come back. Against true widths of 1 to 6 the unbanded rule reports 21, 19, 11, 18, 19 and 12 at \(h = 3\), a slope of \(- 1.09\). The residual rule with independent anchors reports 17, 19, 11, 8, 14 and 13, also \(- 1.09\), with nothing censored at all. A wider range therefore turns a slope of zero into a slope of about \(- 1\): as the true width grows the reported horizon falls. The censored zero understated the failure instead of manufacturing it, and both are reported.

The control on that reading is in the same pass. Every cell whose slope is 1.00 has censoring zero, and rerunning those cells over the wider range returns the identical six estimates, 1 through 6, at every homophily setting. Widening the range moves the censored cells and leaves the uncensored ones untouched, which is what it should do if censoring is what the column says it is.

\textbf{Homophily is one word for several mechanisms, so the grid is rerun under three of them.} A power law on latent distance is one parametrisation, and a result about it is a result about it. We therefore reran the entire grid under two further kernels that are not reparametrisations of the first. One is a bounded logistic kernel, which saturates instead of diverging as distance goes to zero. The other is a \textbf{block} kernel that discards latent distance altogether for discrete latent communities, where two entities are similar if they share a community and not otherwise. That last is a different generative story rather than a different exponent. All three reduce to a uniform weight at \(h = 0\), which keeps the negative control comparable across them.

Two of the three reproduce the collapse and the third does not, and the third is the informative one. Under the power kernel the unbanded rule collapses between \(h = 2\) and \(h = 3\); under the block kernel between \(h = 1\) and \(h = 2\); under the logistic kernel not at all inside the swept range. That is not kernel dependence. \textbf{The parameter} \(h\) \textbf{is not comparable across kernels}, because each concentrates the edge weight differently at the same value of it: the block kernel collapses earliest, the power kernel next, and the logistic kernel is the mildest of the three at every \(h\). A threshold quoted in \(h\) is therefore a statement about one kernel's parametrisation and about nothing else, and quoting one was a defect in how this result was reported.

\textbf{So the collapse is reported against a quantity a field corpus also has.} The natural one is already what the anchor band is built from: the strongest anchor's \textbf{clique contrast at level 1} on the never-event stratum, read at one period. It is a property of a corpus rather than of a generator, it is what ``the anchors carry the target relation's structure'' means once measured, and both sides of the comparison have it. It has to be that statistic and not the more familiar one: Table 7's 0.2448 for trade flows is a matched \emph{density} contrast, and comparing it with a clique contrast would be the units error Section 6.4 documents.

On that scale the two kernels that reach the collapse agree about where it is. The unbanded rule loses the truth at an anchor clique contrast of \textbf{0.0152} under the power kernel and \textbf{0.0186} under the block kernel, a difference of about a fifth between two mechanisms that share no functional form. The logistic kernel does not contradict them: it never exceeds \textbf{0.0063} at any \(h\) we swept, which is below both thresholds, so it is silent rather than contrary. It does not resist the collapse, it fails to reach the concentration at which the collapse happens.

\textbf{And corpus A is inside that regime, measured rather than assumed.} Its strongest anchor, trade flows, has a clique contrast of \textbf{0.0259} at level 1 and one period, against thresholds of 0.0152 and 0.0186. The corpus therefore sits above both, by a factor of 1.7 and 1.4. This is the closest thing to a bridge between the generator and the field that we can build, and it is worth being exact about what it does and does not carry. It does not show that corpus A's phantom horizon is caused by homophily; the trace there is unknown and Section 6.5 says why it must stay unknown. What it does show is that the regime in which the pooled statistic provably loses identification is not an extrapolation past anything real: corpus A's own anchors are more concentrated than the point at which two different mechanisms both break the statistic. Section 10 previously recorded as a threat that the generator could not show any real corpus reaching the level at which identification fails. On this scale it can, and the threat is narrowed rather than removed, because one corpus above a threshold measured on a generator is a weaker thing than a field measurement of the trace itself.

\begingroup\small
\begin{longtable}{@{}
  >{\raggedright\arraybackslash}p{(\linewidth - 6\tabcolsep) * \real{0.1703}}
  >{\raggedright\arraybackslash}p{(\linewidth - 6\tabcolsep) * \real{0.2154}}
  >{\raggedright\arraybackslash}p{(\linewidth - 6\tabcolsep) * \real{0.2222}}
  >{\raggedright\arraybackslash}p{(\linewidth - 6\tabcolsep) * \real{0.3920}}@{}}
\caption{The four passes that answer what one seed, one range and one kernel cannot. The first reruns the entire grid at three seeds and reports the largest per-seed spread of the slope any homophily setting produces.}\label{tab:5}\\
\toprule
\begin{minipage}[b]{\linewidth}\raggedright
\textbf{Pass}
\end{minipage} & \begin{minipage}[b]{\linewidth}\raggedright
\textbf{Cell}
\end{minipage} & \begin{minipage}[b]{\linewidth}\raggedright
\textbf{Estimator}
\end{minipage} & \begin{minipage}[b]{\linewidth}\raggedright
\textbf{Result}
\end{minipage} \\
\midrule
\endfirsthead
\toprule
\begin{minipage}[b]{\linewidth}\raggedright
\textbf{Pass}
\end{minipage} & \begin{minipage}[b]{\linewidth}\raggedright
\textbf{Cell}
\end{minipage} & \begin{minipage}[b]{\linewidth}\raggedright
\textbf{Estimator}
\end{minipage} & \begin{minipage}[b]{\linewidth}\raggedright
\textbf{Result}
\end{minipage} \\
\midrule
\endhead
\bottomrule
\endlastfoot
three kernels & whole grid, power kernel & unbanded rule & collapses between \(h = 2\) and 3, at an anchor clique contrast of 0.0152 \\
three kernels & whole grid, block kernel & unbanded rule & collapses between \(h = 1\) and 2, at an anchor clique contrast of 0.0186 \\
three kernels & whole grid, logistic kernel & unbanded rule & never inside the range, and never exceeds an anchor contrast of 0.0063, below both thresholds \\
three kernels & whole grid, all three & focal isolation, clique estimate & slope 1.00 throughout, in both anchor regimes \\
reference & corpus A, strongest anchor & none, this is the field value & anchor clique contrast 0.0259 at level 1 and one period, above both thresholds \\
three seeds, whole grid & all 36 settings, shared anchors & focal isolation, clique estimate, residual rule & slope 1.00 at every setting and every seed; largest spread 0.000 \\
three seeds, whole grid & all 36 settings, independent anchors & focal isolation, clique estimate & slope 1.00 at every setting and every seed; largest spread 0.000 \\
three seeds, whole grid & independent anchors & residual rule & largest spread 0.273, at \(h = 3\) \\
three seeds, whole grid & both regimes & unbanded rule & largest spread 0.786, at \(h = 2\), which is the transition \\
five seeds & \(h = 1\), both regimes & all four & slope 1.00, spread 0.00, every seed exact \\
five seeds & \(h = 2\), shared anchors & residual rule & slope 1.00, spread 0.00, every seed exact \\
five seeds & \(h = 2\), independent & residual rule & slope 0.88, range 0.54 to 1.00, SD 0.20; mean error +0.43, range 0.00 to +2.00 \\
five seeds & \(h = 2\), both regimes & unbanded rule & slope 0.87, range -0.03 to +1.51, SD 0.56 \\
five seeds & \(h = 3\), both regimes & clique estimate & slope 0.97, range 0.86 to 1.00, SD 0.06; mean error +0.03 \\
five seeds & \(h = 3\), independent & residual rule & slope 0.14, range -0.03 to +0.46, SD 0.22; mean error +4.67 \\
five seeds & \(h = 3\), both regimes & unbanded rule & slope 0.08, range -0.03 to +0.43, SD 0.20; mean error +5.30 \\
wider range & \(h = 3\), independent & residual rule & 9,9,9,8,9,9 to 20: 17,19,11,8,14,13; slope -0.03 to -1.09; censored 6 to 0 \\
wider range & \(h = 3\), both regimes & unbanded rule & 9,9,9,9,9,9 to 20: 21,19,11,18,19,12; slope 0.00 to -1.09; censored 6 to 1 \\
wider range & \(h = 4\), independent & residual rule & 9,9,9,9,9,9 to 20: 19,21,21,17,19,16; slope 0.00 to -0.71; censored 6 to 2 \\
wider range & \(h = 4\), both regimes & unbanded rule & 9,9,9,9,9,9 to 20: 21,21,21,17,21,16; slope 0.00 to -0.83; censored 6 to 4 \\
wider range & every uncensored cell & clique estimate, shared residual & 1,2,3,4,5,6 at both ranges; unchanged \\
\end{longtable}
\endgroup

\textbf{Three seeds across the whole grid, and five where the collapse happens.} One corpus per cell is not an interval, and a reader is entitled to ask about the cell rather than the shape. The whole-grid pass answers it for the columns the argument rests on: over three seeds at all 36 settings of both regimes, focal isolation and the clique estimate return a slope of exactly 1.00 every time, with a largest spread of 0.000. The seed variance in this table is confined to the cells that are collapsing, which is where variance belongs: 0.786 for the unbanded rule at its transition and 0.273 for the residual rule at \(h = 3\) with independent anchors. The five-seed pass on the transition rows is the stricter one and it finds the one exception, below. Table 5's upper block reruns the rows where an estimator is neither exact nor fully collapsed, five seeds each, and it moves two entries and confirms the rest. The entry that moves most is the residual rule at \(h = 2\) with independent anchors, which Table 4 reports as \(+ 2.00/ + 0.54\). Over five seeds the mean slope is \(0.88\) with a range of \(0.54\) to \(1.00\). The single-seed cell is therefore the least favourable of the five, and the transitional character of the row survives while its severity does not. The second is the clique estimate at \(h = 3\), which Table 4 reports as exactly \(1.00\) in both regimes: over five seeds it is \(0.97\) with a range of \(0.86\) to \(1.00\). That is a small correction and it is the right direction to report it in, because ``exact everywhere'' was a claim about one draw. Nothing in the upper block reverses a sign or moves an estimator across the boundary between tracking the truth and not tracking it.

\textbf{The band repairs it only when the anchors share the target relation's structure.} With shared anchors the residual rule is exact at every homophily setting: the band subtracts the term it was designed to subtract. With independent anchors it degrades one step of homophily behind the unbanded rule, slope \(+ 0.54\) at \(h = 2\) and \(- 0.03\) at \(h = 3\). \textbf{Shared structure is what makes the band work and independent structure is what breaks it}, which is the opposite of the direction the mechanism is usually stated in and the opposite of the direction we stated it in. Corpus B is the independent case, and the reason its residual of 0.4075 does not become a phantom horizon is that it has no time axis for a horizon to live on. Mechanism (h) is stated accordingly in Section 7.

\textbf{Focal isolation is unbiased in homophily everywhere.} The first column is 0.00 / 1.00 in all twelve cells, six homophily settings in each of the two anchor regimes. It reads the group's own record, and no amount of similarity between the candidate and the members puts the candidate on that record. This is the quantitative form of the argument Sections 6.4 and 6.5 make across two field sources, and it is made here on one corpus against one known truth.

\textbf{Saturation and homophily are one mechanism, not two.} The clique estimate's column is exact everywhere in Table 4 and very nearly exact over five seeds, which is a real result and a narrow one. A clique needs every dyad, so in a relation at density 0.12 homophily cannot complete one: the statistic is specific for a combinatorial reason. Raising the target relation's density at fixed homophily removes that protection. At \(h = 3\) and width \(\beta = 2\), the estimate is exact at density 0.12, does not fire at all inside eight periods at density 0.50, and overshoots by two at density 0.90; at \(h = 0\) it is exact at all three densities. Corpus A's target relation has density 0.9954, and at its declared primary level the clique estimate does not fire anywhere inside ten periods. Mechanisms (b) and (h) are therefore not independent findings about two corpora: saturation is the condition under which homophily breaks the clique statistic, and the manuscript's own separation of them understated the interaction.

\textbf{A heterogeneous width makes the estimand a choice.} Redraw the width per event from the PTA register's coded gaps. The mean presence at lead time \(\Delta\) is then the share of events whose width is at least \(\Delta\), so a rule that waits for that mean to reach zero returns the maximum width and a rule that thresholds it at one half returns the median. Both are exact for their own estimand in all twelve settings: the zero rule returns the realised maximum plus one and the median rule returns the realised median plus one, every time. We report the maximum-based rule, for the reason a step function's mean is a width rather than a presence: one non-zero lead time is enough to widen the horizon. The pre-specified estimate also returns the maximum plus one in all twelve settings, while the two pooled rules wander between three and nine. A maximum-based rule is decided by the single widest event in the corpus, which is worth saying beside the exactness: it is why the rule is conservative and it is also why one badly coded event moves it.

\textbf{Dating noise reproduces the declared-resolution stratum.} Displacing a share of coded dates one period late makes the register wrong rather than early. That puts focal isolation's presence at exactly that share at one period and at zero beyond: 0.0667 at a displaced share of a tenth, 0.3000 at three tenths. Corpus A's declared-resolution stratum reads 0.1419 at one period, which under this model corresponds to about one event in seven in that stratum carrying a coded date one period late. That is a quantitative bridge between the synthetic testbed and the one place in corpus A where back-dating is measured rather than entailed.

\subsection{Where the Administrative Trace Lives}

Section 3.3 states the claim this section tests, and the exact objection an outside reading raised against it. Administrative back-dating of the focal event can reach \(r^{\star}\) through a neighbouring group's record, which no reading of the focal register opens. It correlates with the label when the two entries fall near each other, which is what an accession wave is. Table 4's first column presents focal isolation as unbiased on the strength of that claim, so the objection is load-bearing rather than a detail.

The generator of Section 5.2 was extended to make the wave a controlled parameter, and it records per event which record carries the trace and how far it reaches, so the question is decidable rather than arguable. We sweep the wave share against six widths, four homophily settings and five seeds in each of two convention regimes, 840 corpora. The edge budget is sized from the measured worst case over the whole sweep, which a wave enlarges: the widest setting peaks at 2,427 eligible dyads of 19,900 pairs. The density is therefore 0.14 against Table 4's 0.12, chosen as the first covering value of a declared ladder rather than by hand. Within the sweep \(r^{\star}\)'s density is again one exact value across every cell, and the smallest headroom any cell had is 354 dyads.

Three truths are graded against, and keeping them apart is the whole difficulty.

\begin{itemize}
\tightlist
\item
  \textbf{The focal-record horizon} is what a per-group audit reads: the largest width the focal group's own register carries, plus one. This is what Section 11's first instruction told a practitioner to measure.
\item
  \textbf{The every-register horizon} is what the pipeline actually computes, because focal isolation is run over every group's own record and not only over the group whose event is being audited.
\item
  \textbf{The administrative horizon} is the truth: the largest lead time at which any record carries a dyad from the acceding entity into the focal group's members that the honest registers do not, counting only dyads the focal event caused.
\end{itemize}

The last restriction is not a convenience. Any group's back-dating of any membership puts administrative dyads into \(r^{\star}\) at arbitrary lead times, and an entity that joined some other group nine periods ago carries one at lead eleven for reasons that have nothing to do with the focal event. Calling that a leakage horizon of eleven would be the error this paper is about. That ambient content is reported separately and against matched negatives, below.

\textbf{Under one dating convention the claim holds.} In all 480 corpora of that regime, at wave shares from zero to half the events and with 8 to 45 partner entries added per corpus, \textbf{not one event} carried an off-focal administrative trace that the focal record did not already carry. The gap between the focal-record horizon and the administrative horizon is exactly zero in every cell. Focal isolation is exact with slope 1.00 at every wave share, and so is the residual rule; the clique estimate is exact up to a quarter of events and reads +0.01 with slope 0.99 at a half, which is one corpus of 120 moving by one period. That is Table 4's reading reproduced on corpora containing the mechanism that was supposed to break it. The reason is simple once measured: the partner register is back-dated by the same width as the focal one, so every dyad the wave puts into \(r^{\star}\) early is one the focal record already shows.

\textbf{Under mixed conventions the claim fails, and by up to nine periods.} Give the partner register the trade register's own coded gap distribution, whose tail reaches nine periods, against a sharp focal width. The focal-record reading then undershoots the administrative horizon by a mean of 0.82 periods at a wave share of a tenth, 2.22 at a quarter and 2.68 at a half, with a maximum of nine and with 44, 65 and 78 percent of cells affected. The share of events carrying a caused off-focal trace wider than their own record shows runs 1.8, 5.1 and 8.9 percent at the three wave shares, and reaches 42 percent in the worst single corpus. This is corpus A's situation rather than a contrivance: that corpus dates 7,615 of its accessions exactly and 447 only to a five-year block, inside one target relation.

\textbf{The failure is administrative and not a form of homophily.} Under mixed conventions the mean gap is 3.77, 2.07, 2.67 and 2.23 periods at homophily 0, 1, 2 and 3, and under one convention it is 0.00 at all four. The mechanism therefore does not switch on with similarity, which is what distinguishes it from mechanism (h) and why Section 3.4 gives it its own term. A reference band that removed homophily perfectly would leave this untouched.

\textbf{Reading every register closes most of it, which is the recommendation.} Figure 3 puts the two readings against the truth in both regimes. Graded against the administrative horizon, the every-register reading errs by \(- 0.81\), \(- 0.87\) and \(- 1.32\) periods at the three wave shares, that is it now exceeds the truth by about a period rather than falling short of it by two or three, with a residual shortfall of at most five periods in the worst cell. Overshooting is the direction Section 3.2's loss prefers. The remedy is procedural and cheap: focal isolation reads a membership register and no graph, so running it over every group costs no more than running it over one.

\textbf{The pooled statistics are not a substitute.} They do not recover what the focal reading misses. This is the result we would have wanted to go the other way. The clique estimate and the residual rule read \(r^{\star}\) and therefore do see the off-focal dyads, so a reader might expect them to bracket the truth from above. They do not: against the administrative horizon the clique estimate undershoots in 39 to 53 of 120 cells with slopes of 0.15 to 0.99, and the residual rule in 53 to 91 of 120 with slopes of 0.12 to 1.19. Seeing a dyad is not the same as attributing it, and a statistic that pools every dyad cannot tell an administrative one from a similar one, which is the same finding arriving from the other side.

\textbf{The ambient content is not a trace, and the control says so.} Off-focal administrative dyads exist in these corpora whether or not the focal event happened, and they reach as far as thirteen periods. Across the whole sweep the mean reach is 1.0053 periods for positives and 0.8396 for matched negatives drawn from the same group at the same period, and the share of candidates carrying one at all is 0.1401 against 0.1359. Those are the same number to within a rounding of the share. An administrative dyad that positives and negatives carry alike is not a trace of the event, and a diagnostic that read it as one would be reporting exactly the artefact this paper is about. That is why the caused part is separated and the ambient part reported beside it rather than added to it.

\begin{figure}[!htbp]
\centering
\includegraphics[width=\textwidth]{figures/figure-3.pdf}
\caption{\textbf{The focal record holds the whole trace under one dating convention and not under two.} Synthetic testbed with a curatorial wave as a controlled parameter, 840 corpora over six widths, four homophily settings and five seeds.}
\end{figure}

\begin{table}[!htbp]
\centering
\small
\begin{tabular}{@{}
  >{\raggedright\arraybackslash}p{(\linewidth - 10\tabcolsep) * \real{0.1788}}
  >{\centering\arraybackslash}p{(\linewidth - 10\tabcolsep) * \real{0.1769}}
  >{\centering\arraybackslash}p{(\linewidth - 10\tabcolsep) * \real{0.1358}}
  >{\centering\arraybackslash}p{(\linewidth - 10\tabcolsep) * \real{0.1536}}
  >{\centering\arraybackslash}p{(\linewidth - 10\tabcolsep) * \real{0.1705}}
  >{\centering\arraybackslash}p{(\linewidth - 10\tabcolsep) * \real{0.1844}}@{}}
\toprule
\begin{minipage}[b]{\linewidth}\raggedright
\textbf{Convention}
\end{minipage} & \begin{minipage}[b]{\linewidth}\centering
\textbf{Wave share}
\end{minipage} & \begin{minipage}[b]{\linewidth}\centering
\textbf{Gap, focal record only}
\end{minipage} & \begin{minipage}[b]{\linewidth}\centering
\textbf{Gap, every register}
\end{minipage} & \begin{minipage}[b]{\linewidth}\centering
\textbf{Cells with a shortfall}
\end{minipage} & \begin{minipage}[b]{\linewidth}\centering
\textbf{Events with an off-focal trace}
\end{minipage} \\
\midrule
one & 0 & 0.00, max 0 & 0.00, max 0 & 0 percent & 0.0 percent \\
one & 0.1 & 0.00, max 0 & 0.00, max 0 & 0 percent & 0.0 percent \\
one & 0.25 & 0.00, max 0 & 0.00, max 0 & 0 percent & 0.0 percent \\
one & 0.5 & 0.00, max 0 & 0.00, max 0 & 0 percent & 0.0 percent \\
mixed & 0.1 & +0.82, max 4 & -0.81, max 1 & 44 percent & 1.8 percent \\
mixed & 0.25 & +2.22, max 9 & -0.87, max 5 & 65 percent & 5.1 percent \\
mixed & 0.5 & +2.68, max 9 & -1.32, max 5 & 78 percent & 8.9 percent \\
\bottomrule
\end{tabular}
\caption{Where the administrative trace lives, by dating convention and wave share. 840 corpora, six widths, four homophily settings, five seeds, at the density Table 4 uses.}\label{tab:6}
\end{table}

Two things this experiment does not settle. It cannot say which regime corpus A is in as a magnitude, only that corpus A has two dating conventions inside one target relation and is therefore not in the regime where the claim holds. And the examined range of eight periods censors focal isolation itself here, because the partner law reaches nine. At wave shares of a quarter and a half the every-register reading loses slope for that reason rather than from any failure of the operation. That is Section 6.2's censoring lesson applying to the instrument we recommend rather than to the one we criticise. The range has to be at least as wide as the widest convention in the corpus, and that is a declaration a practitioner has to make before looking.

\subsection{A Width Printed in the Source}

Section 5.3's register codes the back-dating width per agreement. We ask the same question two ways.

Measured on the focal agreement alone, presence is 1 for every lead time inside the coded gap and 0 for every lead time outside it, at every gap width, with no decline as the gap widens. That result is entailed by the construction, since the relation is dated from signature, the event from entry into force, and the gap is their difference; it verifies the pipeline rather than recovering the width independently, and Table 1 records it as such. The entailment is what makes the comparison sharp rather than what weakens it: the width is in the data by construction, so an instrument that misses it is definitely missing something that is definitely there.

Measured the way the suite specifies, as a density over the whole relation, the contrast between inside and outside the coded gap is \textbf{0.1023}, with a 95 percent interval of 0.0905 to 0.1139 over 758 agreements. The curve is nearly flat: 0.7519 at zero lead time falling only to 0.6206 at eight periods. Agreements with a coded gap of three periods still show 91 percent presence eight periods out. Figure 4 puts the two measurements side by side.

\textbf{The mechanism is a share, and the share is the result.} The signature-dated relation carries 204,747 edge-years against the force-dated relation's 190,781, so the back-dated dyads are 13,966 edge-years, 6.8 percent of the relation. A density taken over the whole relation is therefore a weighted average in which the quantity of interest has weight 0.068 and the acceding party's other agreements have weight 0.932. That is why the pooled curve is flat where the focal curve is a step, and it is the whole explanation.

We state it as a share rather than as a ratio, because the obvious ratio does not survive a look at its units. The focal reading is a \textbf{presence probability}, exactly 1 inside the coded gap by entailment; the pooled reading is a \textbf{density contrast}, 0.1023. Their quotient is 9.77, with an interval of 8.78 to 11.05 coming entirely from the denominator, and it is tempting to call that the pooled statistic recovering a tenth of the width. It is not, because the pooled statistic does not estimate a width at all: a width is a whole number of \emph{periods}, here running from zero to nine with a median of one, and 0.1023 is dimensionless. What 9.77 measures is how far apart two instruments read on one set of events.

The same comparison on corpus A, on the same events and the same lead times, gives a step from zero lead time to the mean beyond it of \textbf{0.9844} under focal isolation against \textbf{0.1447} on the pooled relation, a ratio of \textbf{6.80} with an interval of 4.74 to 10.61. The two corpora therefore give ratios of 9.8 and 6.8, and we report them separately: they are 44 percent apart as point estimates, no mechanism forces the ratio to any particular value, and there are two observations. They are also not distinguishable, because the coalition corpus's interval contains the trade register's point estimate. Neither licenses a claim that the ratio is a constant of the instrument, which is what calling the two readings one factor would license. What both corpora do support is weaker and more useful: the group's own record and the pooled relation disagree by close to an order of magnitude on the same events, and the size of the disagreement is corpus specific. Both ratios are between two statistics on different scales, so both carry the caveat above; they are reported because the disagreement is the finding, not because either quotient is a width.

\begin{figure}[!htbp]
\centering
\includegraphics[width=\textwidth]{figures/figure-4.pdf}
\caption{\textbf{A width printed in the source, and two instruments that part company on it.} The trade agreement register of Section 5.3, 758 agreements and 6,822 event-lead-time pairs, presence at level 1 on the signature-dated relation, lead times 0 to 8 measured back from the coded entry-into-force year.}
\end{figure}

\subsection{The Coalition Corpus}

We ran the suite as specified on Section 5.4's corpus, then ran focal isolation on the same events. Before the numbers, what they can and cannot establish, because the two are easy to conflate and the difference decides how much of this section carries weight.

\textbf{What this corpus cannot do.} It cannot falsify the pooled reading by comparison with focal isolation. The register codes one date per membership for 94 percent of its events, so focal isolation's zero there follows from the register's structure and not from a measurement, and a statistic that is silent by construction cannot contradict one that speaks. That much Section 3.4 and Table 1 already state. What is new here is that the obvious repair does not exist either. The one stratum where focal isolation measures is the declared-resolution stratum, and the register assigns that stratum only to organisations founded before 1965, so all 447 of its accessions carry a coded year of 1965 or earlier while the analysis window opens in 1965. \textbf{The deepest lead time at which the pooled statistic could be read on those events, inside the window, is zero.} Reading it at one to eight periods would mean reading 1957 to 1964, inside the coverage break the window exists to exclude, and 86 percent of the rows such a reading would use fall there. So on corpus A the two instruments cannot both be measurements on the same events, anywhere, and no subsetting repairs it. Where focal isolation measures, the pooled statistic is not computable without the confound the design removed; where the pooled statistic is computable, focal isolation is entailed.

\textbf{What this corpus does do, without any of that.} Three of the four results below need no focal isolation at all, and they are what the section rests on. The target relation is statistically indistinguishable at eight periods from a relation that cannot record the event, which is a falsification of the pooled reading on its own terms. The lead time at which the pre-specified rule declares the horizon closed ranges over the whole examined interval under resampling, which says the width is not identified whatever it is. And the four-period decay the unconditioned curve shows is almost entirely prior membership spells, which is a competing explanation the pooled statistic cannot distinguish from a trace. Where the pooled reading is falsified against a \emph{known} width, it is falsified on the trade register of Section 5.3, where both dates are coded, and there \(n\) is 759 agreements. The comparison between the two instruments on corpus A is reported below as a description of how far apart they read, which is what it is.

\textbf{Under focal isolation the corpus carries no trace beyond one period.} Conditioned on first-time entry, the trace is present at zero lead time and exactly absent at every lead time of one or more. It runs 1.000 to 0.000 for the 7,146 exactly dated accessions counted at zero lead time, and 1.000 to 0.000 for the 202 multilateral treaty accessions. One stratum does carry a trace at one period, the declared-resolution stratum, and it is set out three paragraphs below because it is the only measurement of the three readings in this section. The back-dating is therefore at most one period, and zero for 94 percent of the corpus.

\textbf{How much of that zero is measured and how much is entailed.} Corpus A's register codes \textbf{one} date per membership and derives the accession event from a transition in the same field. So for the exactly dated stratum a zero residual follows from the register's structure rather than from a measurement, in the same way that Section 6.4's step table follows from its construction. Section 3.4 states the condition: focal isolation identifies a width only where the source dates the relation and the event from separate fields. This register does not, for 94 percent of its events, and no reading of that zero says the true horizon is zero. What the comparison establishes is not ``the true horizon is zero and the diagnostic says eight''. It is that \textbf{the focal group's own record carries nothing past one period while the pooled statistic reports a contrast still at 0.1490 after eight}. That gap is the diagnostic's blind spot measured on a corpus of realistic size, rather than a demonstration that the pooled number is false. On the step from zero lead time to the mean beyond it the two instruments give 0.9844, with an interval of 0.9809 to 0.9878, and 0.1447, with an interval of 0.0928 to 0.2075. Section 6.4 supplies the case where the width is independently coded and the pooled reading is therefore demonstrably wrong; this section supplies the size of the gap at scale.

\textbf{The one place back-dating is measured rather than entailed.} The declared resolution stratum reads \textbf{0.1419} at one period and 0.000 beyond, over 310 events in 42 families, with a 95 percent interval of 0.0678 to 0.2395. The exactly dated stratum reads 0.0 at the same lead time with an interval of zero width, which is what an entailment looks like beside a measurement. That reading is over \textbf{7,084} events rather than the 7,146 of the paragraph above, and the difference is the lead time the count is taken at rather than a different stratum. Focal isolation is defined at lead time one only where the focal group has a member set in the preceding period. 62 of the exactly dated first entries are accessions in a group's founding period, where there is no earlier record for the entity to be absent from. Section 5.5's flow states that filter with the others. This is not an entailment and it is the paper's only positive field measurement of a leakage horizon on the coalition corpus, so it is worth stating precisely what produced it. The register declares that its own dates before 1965 are accurate only within five year blocks, and the coding rule that locates an event inside such a block is exactly the kind of curatorial decision Section 1 describes: a date assigned by procedure rather than observed. One event in seven in that stratum consequently puts the acceding state on the record a period before its coded accession, and the interval admits anything between one in fifteen and one in four. The width is one period; the share of events carrying it is not pinned down by 310 events over 42 families. Section 6.2's dating-noise sweep reproduces that figure from the other direction, displacing a declared share of coded dates by one period and recovering a presence of exactly that share at one period and zero beyond. The width is small, it is one period, and it is real; a diagnostic that reported eight would be wrong by a factor of eight rather than by an unknown amount.

\textbf{The unconditioned curve says four periods, and it is recurrence.} Pooled over all accessions, the declared resolution stratum decays 0.3490, 0.2489, 0.1175, 0.0448, which is the textbook shape of a four-period leakage horizon with a floor. It is prior membership spells, which Section 2 places against the dynamic link prediction literature that established the confound \citep{poursafaei}. What is specific to a leakage diagnostic is that the repetition manufactures a decaying curve with the shape of a horizon rather than merely inflating a score, and that it loads onto whichever stratum recurs most. That stratum's recurrence rate is 30.65 percent against the exactly dated stratum's 6.16 percent, so the artefact loads onto precisely the stratum that appears to have a horizon, which is also the stratum whose dates are least precise. Figure 5 shows the same events before and after conditioning. Any trace statistic must therefore be computed conditional on first entry, and a corpus's recurrence rate must be reported beside any horizon it produces.

\begin{figure}[!htbp]
\centering
\includegraphics[width=\textwidth]{figures/figure-5.pdf}
\caption{\textbf{A four-period horizon that is entirely recurrence.} Corpus A, focal isolation on the membership register, 447 declared-resolution and 7,615 exactly dated accession events, lead times 0 to 8.}
\end{figure}

\textbf{Anchors decay like the target relation.} Table 7 reports the matched contrast by relation and lead time, on 52,305 positives with negatives drawn from the same group and period, and Figure 6 plots the full curves. The statistic is the difference in the \textbf{mean share of the candidate's dyads into the existing members that the relation carries}, positives minus matched negatives, at level 1. It is not a clique rate, and the two are different objects here. The clique contrast appears in Table 3, Table 8 and Figure 2; this one is a density.

\begin{table}[!htbp]
\centering
\small
\begin{tabular}{@{}
  >{\raggedright\arraybackslash}p{(\linewidth - 10\tabcolsep) * \real{0.2105}}
  >{\centering\arraybackslash}p{(\linewidth - 10\tabcolsep) * \real{0.1118}}
  >{\centering\arraybackslash}p{(\linewidth - 10\tabcolsep) * \real{0.1745}}
  >{\centering\arraybackslash}p{(\linewidth - 10\tabcolsep) * \real{0.1745}}
  >{\centering\arraybackslash}p{(\linewidth - 10\tabcolsep) * \real{0.1745}}
  >{\centering\arraybackslash}p{(\linewidth - 10\tabcolsep) * \real{0.1541}}@{}}
\toprule
\begin{minipage}[b]{\linewidth}\raggedright
\textbf{Relation}
\end{minipage} & \begin{minipage}[b]{\linewidth}\centering
\textbf{Role}
\end{minipage} & \begin{minipage}[b]{\linewidth}\centering
\[\mathbf{\Delta =}\mathbf{0}\]
\end{minipage} & \begin{minipage}[b]{\linewidth}\centering
\[\mathbf{\Delta =}\mathbf{4}\]
\end{minipage} & \begin{minipage}[b]{\linewidth}\centering
\[\mathbf{\Delta =}\mathbf{8}\]
\end{minipage} & \begin{minipage}[b]{\linewidth}\centering
\textbf{Retention at 8}
\end{minipage} \\
\midrule
Co-membership & target & 0.3114 {[}0.2504, 0.3773{]} & 0.2038 {[}0.1725, 0.2314{]} & 0.1490 {[}0.1160, 0.1768{]} & 47.9 percent \\
Trade flows & anchor & 0.2639 {[}0.2415, 0.2825{]} & 0.1890 {[}0.1604, 0.2168{]} & 0.1463 {[}0.1161, 0.1755{]} & 55.4 percent \\
Voting agreement & anchor & 0.1779 {[}0.1549, 0.2036{]} & 0.1296 {[}0.1053, 0.1581{]} & 0.1081 {[}0.0842, 0.1376{]} & 60.8 percent \\
Alliance ties & other & 0.1179 {[}0.0905, 0.1521{]} & 0.0850 {[}0.0632, 0.1133{]} & 0.0670 {[}0.0482, 0.0901{]} & 56.8 percent \\
Militarised disputes & anchor & 0.0041 {[}0.0032, 0.0051{]} & 0.0038 {[}0.0028, 0.0053{]} & 0.0044 {[}0.0032, 0.0061{]} & flat \\
\bottomrule
\end{tabular}
\caption{Matched \textbf{density} contrast by relation and lead time, corpus A, with 95 percent intervals from 1,000 bootstrap replicates resampling the 464 membership equivalence families. 52,305 positives over 303 families that carry at least one event.}\label{tab:7}
\end{table}

Trade flows cannot record an accession to an intergovernmental organisation, and they produce a monotone decaying contrast retaining 55.4 percent of its value at eight periods. So does voting agreement, at 60.8 percent. The one relation that comes out flat is militarised disputes, at 0.004 everywhere, and it is the one relation with no homophily content. What every non flat series measures is that entities which accede to a group are genuinely more similar to its members than never joiners are, and were so a decade earlier. That is the signal a forecaster should be using, not a trace of the event.

The intervals sharpen the point past what the point estimates can carry. At eight periods the target relation's interval is 0.1160 to 0.1768 and trade flows' is 0.1161 to 0.1755. \textbf{The relation built from the label source and a relation that cannot record the label are indistinguishable at eight periods before any band is applied.} At zero lead time they are separated, 0.2504 to 0.3773 against 0.2415 to 0.2825, so the two series do part company where the event is. They have stopped doing so by the far end of the range the diagnostic examines.

\textbf{A negative control, and the only relation the instrument separates from the target.} Militarised disputes sit at 0.004 at every lead time with an interval two thousandths wide, and the interval at eight periods overlaps the interval at zero: no decay at all. This is the paper's cleanest evidence that the instrument is not simply blind, and it is the second of the three results in this section that needs no focal isolation. The one relation that can record neither the event nor the similarity between the acceding state and the members is the one relation the diagnostic separates from the target, and it is the relation a documentary reading would have nominated in advance. \textbf{This is the field half of a pair.} Its synthetic half is Section 6.1's \(h = 0\), \(\beta = 0\) cell, where there is nothing to find and every estimator returns the correct value of 1. One is a negative control on a corpus whose truth is unknown and the other a negative control where the truth is set, and together they say the instrument's zeros are readable: it is not reporting a horizon everywhere.

\textbf{The band works, and its crossing lead time is not identified.} The target relation minus the strongest anchor runs 0.0475, 0.0667, 0.0327, 0.0252, 0.0148, 0.0041, first falling below 0.01 at \(\Delta = 5\), so the pre-specified rule returns about five periods and focal isolation on the same events says zero. That five is a point estimate and the resampling does not support it as a measurement. Across 1,000 family-level replicates the crossing lead time has a median of 4 and a 95 percent interval spanning 0 to 7. It never crossed inside eight periods in 222 replicates, and in 133 replicates the residual was already below the threshold at zero lead time, which is to say there was no phantom at all. The modal value is 5, at 261 replicates, and the next is 4, at 201. \textbf{The phantom is real in most samples and its width is not identified.} That is a more useful statement than ``about five periods'' and a more damaging one for the instrument: a diagnostic whose output ranges over the whole examined interval under resampling is not measuring a width.

This is a third distinct reason for non-identification and the deepest. The first is noise: on the 59 alliance entities where the statistic has power in principle, 177 events give a standard error equal to the effect. The second is saturation, and it has to be quantified on this corpus rather than borrowed from Section 5.3. Of the 474,034 candidate-to-member dyads the pooled statistic reads as present at \(k = 1\), \textbf{117} are present only through the focal entity: 0.02 percent. At the declared primary level \(k = 25\) the figure is 15,059 of 230,707, or \textbf{6.5 percent}. Co-membership carries a weight equal to the number of organisations a pair shares, so the decomposition is exact: a dyad read as present at level \(k\) is present only through the focal entity exactly when its weight equals \(k\). The level therefore decides the answer by a factor of 260, and at the level a reader would assume the focal group has contributed one dyad in four thousand of what the statistic reads. The third reason is neither noise nor saturation, and no amount of data or thresholding touches it, because a trace of the event and a persistent similarity to the members enter the statistic identically. Section 6.2 measures what that costs: with no band at all the reported horizon becomes independent of the true width once homophily is strong enough.

\textbf{Where the generator's settings meet this corpus.} Section 10 states, correctly, that a generator can show homophily destroying identification and cannot show that any real corpus reaches the strength at which it does. Table 7 is the empirical half of that argument, and it is stronger than a possibility proof. Trade flows cannot record an accession to an organisation and their contrast retains 55.4 percent of its value at eight periods against the target relation's 47.9 percent; voting agreement retains 60.8 percent. Anchors that decay like the target relation are what ``the anchors share the target relation's latent structure'' means when it is measured rather than assumed, and it is the condition Section 6.2's shared-anchor regime sets by hand. The generator's high-homophily cells are therefore not an arbitrary extrapolation from this corpus: they are the regime this corpus's own anchor curves put it in.

\begin{figure}[!htbp]
\centering
\includegraphics[width=\textwidth]{figures/figure-6.pdf}
\caption{\textbf{Relations that cannot record the event decay like the one that can.} Corpus A, 1965 to 2014, matched star density per relation at level 1, 52,305 positives over 303 membership equivalence families with negatives drawn from the same group and period, lead times 0 to 8.}
\end{figure}

\subsection{A Scope Test on the Inverse Size Distribution}

Three of the mechanisms in Section 7 are claims about any curated group corpus, and a fourth, saturation, is a claim whose scope a second domain can settle. We test all four on corpus B, whose size distribution is the inverse. Statistics are computed by the same pinned implementation used for corpus A, so a difference between domains cannot be an implementation difference. Corpus B has no time axis, so it cannot reproduce a single result about lead time, and lead time is what a leakage horizon is measured in. Calling it a replication would promise otherwise.

\textbf{The size dependence of the clique statistic generalises.} Table 8 reports the contrast by group size for both corpora, computed by one script so the comparison is not across implementations. Both decay with size and both are identically zero above eighty members, in domains sharing no entities, no relations and, for corpus B, no time axis. The two sample opposite ends: 5,486 of corpus A's 6,597 events fall in groups above twenty members, and 1,399 of corpus B's 2,540 complexes at three members or fewer.

The two series share a shape and not a level. They sit at visibly different heights in every bin where they overlap, and they are read at different levels of the density ladder, so we do not compare the heights. Corpus A is read at its declared primary level \(k = 25\), because at \(k = 1\) its weighted target relation has density 0.9954 and the statistic is degenerate there. Corpus B is read at \(k = 1\), because its channels are binary and no other level exists. Both are taken at zero lead time, since corpus B has no other. The claim the second domain supports is therefore this: \textbf{in both domains the clique contrast decreases monotonically with group size and is structurally zero above eighty members, and the two corpora sample opposite ends of the same range.} They are two decreasing functions read at two levels rather than one function traced twice, which is what Figure 7 draws.

\begin{table}[!htbp]
\centering
\small
\begin{tabular}{@{}
  >{\raggedright\arraybackslash}p{(\linewidth - 12\tabcolsep) * \real{0.1684}}
  >{\centering\arraybackslash}p{(\linewidth - 12\tabcolsep) * \real{0.1308}}
  >{\centering\arraybackslash}p{(\linewidth - 12\tabcolsep) * \real{0.1424}}
  >{\centering\arraybackslash}p{(\linewidth - 12\tabcolsep) * \real{0.1424}}
  >{\centering\arraybackslash}p{(\linewidth - 12\tabcolsep) * \real{0.1541}}
  >{\centering\arraybackslash}p{(\linewidth - 12\tabcolsep) * \real{0.1308}}
  >{\centering\arraybackslash}p{(\linewidth - 12\tabcolsep) * \real{0.1308}}@{}}
\toprule
\begin{minipage}[b]{\linewidth}\raggedright
\textbf{Members}
\end{minipage} & \begin{minipage}[b]{\linewidth}\centering
\textbf{2 to 3}
\end{minipage} & \begin{minipage}[b]{\linewidth}\centering
\textbf{4 to 5}
\end{minipage} & \begin{minipage}[b]{\linewidth}\centering
\textbf{6 to 10}
\end{minipage} & \begin{minipage}[b]{\linewidth}\centering
\textbf{11 to 20}
\end{minipage} & \begin{minipage}[b]{\linewidth}\centering
\textbf{21 to 80}
\end{minipage} & \begin{minipage}[b]{\linewidth}\centering
\textbf{81 and above}
\end{minipage} \\
\midrule
Corpus B positives & 1,399 & 641 & 321 & 142 & 34 & 3 \\
Corpus B clique rate, positive & 0.4124 & 0.4665 & 0.4455 & 0.3662 & 0.1765 & 0.0000 \\
Corpus B negative cliques, count & 0 of 1,399 & 0 of 641 & 0 of 321 & 0 of 142 & 0 of 34 & 0 of 3 \\
Corpus B expected negative cliques, independent sampling & \(2.1 \times 10^{- 6}\) & \(1.0 \times 10^{- 15}\) & \(7.5 \times 10^{- 43}\) & \(5.4 \times 10^{- 161}\) & \(0\) & \(0\) \\
Corpus B contrast, \(k = 1\) & 0.4124 {[}0.3855, 0.4386{]} & 0.4665 {[}0.4297, 0.5066{]} & 0.4455 {[}0.3941, 0.4982{]} & 0.3662 {[}0.2889, 0.4516{]} & 0.1765 {[}0.0571, 0.3104{]} & 0.0000 \\
Corpus A positives & 24 & 47 & 289 & 751 & 2,801 & 2,685 \\
Corpus A contrast, \(k = 25\) & 0.5833 {[}0.3148, 0.8096{]} & 0.3617 {[}0.2211, 0.5083{]} & 0.3391 {[}0.2451, 0.4322{]} & 0.2437 {[}0.1781, 0.3151{]} & 0.0521 {[}0.0323, 0.0755{]} & 0.0000 \\
Corpus A contrast, \(k = 1\) & 0.0625 {[}0.0000, 0.1250{]} & 0.1809 {[}0.0980, 0.2675{]} & 0.1972 {[}0.1597, 0.2372{]} & 0.1911 {[}0.1576, 0.2207{]} & 0.2387 {[}0.2018, 0.2740{]} & 0.4778 {[}0.3622, 0.5754{]} \\
\bottomrule
\end{tabular}
\caption{Clique contrast by group size, both corpora, matched or random negatives respectively, with 95 percent intervals from 1,000 bootstrap replicates. Corpus A resamples membership equivalence families; corpus B resamples complexes, which keeps a positive and the negatives generated from it in the same replicate.}\label{tab:8}
\end{table}

\textbf{A zero negative clique rate is what a correct sampler produces.} A clique on \(n\) members needs all \(\binom{n}{2}\) dyads, and the curated-register channel has density 0.00115. Summing \(d^{\binom{n_{i}}{2}}\) over the realised sizes of the smallest bin's 1,399 negatives therefore gives an expected count of two millionths of a group against an observed zero. The zero therefore carries no information about the sampler, and Section 5.5's check that the negatives are not edge-excluded is the statement that does. The stricter comparison is the hard stratum, which keeps most of a real complex and swaps two members out. There the smallest bin yields \textbf{9 cliques of 2,798 negatives} at a contrast of 0.4092 against the random stratum's 0.4124, so the separation is not an artefact of negatives that could not have been cliques.

\textbf{The small bins are small, and their intervals say so.} Corpus A's largest entry, 0.5833 at two to three members, rests on 24 events with an interval of 0.3148 to 0.8096 that overlaps the next bin's, so the apparent peak at the smallest sizes is not resolved. That is mechanism (d) measured: the subpopulation where the statistic has power is the one with too few events. Corpus B's smallest bins carry 1,399 and 641 positives with intervals a twentieth as wide, which is the same mechanism seen from the side where it does not bind.

The last row is the one to read against the two above it. At a level where the target relation is saturated the size dependence \textbf{reverses}, rising with group size instead of falling, because both classes are near-cliques and the contrast then comes from the negative rate declining with size. Section 7 states the consequence for mechanisms (b) and (c).

\textbf{One claim is falsified.} The claim was that group size and sample size trade off in any curated group corpus, so that no subpopulation offers both. Corpus B offers both: 1,399 positives at two to three members, a contrast of 0.4124, and a negative clique rate of 0.0000. The trade off is a property of a corpus whose group size correlates with entity type when the small-size type is rare, which is corpus A exactly, with 59 alliances at median size 5 carrying 177 events against numerous large organisations. The restatement is the more useful claim because it says when to expect the problem.

\textbf{One claim narrows, and the controlled test reverses its explanation.} The finding that the anchors carry the same homophily as the target relation does not transfer. On corpus B the target relation contrasts 0.4240 against a strongest anchor of 0.0165, a residual of 0.4075, and 0.4059 on hard negatives; the four remaining anchors sit at 0.0067, 0.0051, 0.0000 and 0.0000, which is flat. In international relations, trade, voting agreement and co-membership express one latent quantity, so subtracting one from another removes almost nothing; messenger RNA co-expression and gene fusion bear no such relation to a curator's complex boundaries.

Corpus B is therefore the independent-anchor case of Section 6.2, where the band leaves the homophily in rather than subtracting it, and the only reason its residual of 0.4075 is not a phantom horizon is that it has no time axis for one to live on.

\textbf{One claim is confirmed in scope.} No relation in corpus B exceeds density 0.001887 and the target relation is 0.001153 against corpus A's 0.9954, a factor of 864, so nothing here saturates and the saturation mechanism is correctly a statement about coalition relations. Densities are taken over deduplicated unordered pairs against the shared universe of 19,486 proteins, which matters because the source ships both directions of every interaction.

\textbf{One result concerns method rather than instrument.} A limitation travels with it. The three channels declared contaminated on documentary grounds came out ranked first, second and third on clique contrast, at 0.4579, 0.4240 and 0.3075, above every declared anchor at 0.0165 or below. The text mining channel scores above the target relation itself. Reading what each channel records identified all three leaking channels, with no false positives among five anchors. Section 3.2 asserts that no statistic separates leakage from signal and that the verdict requires reading the coding rules; here the reading was written down and the statistic then placed exactly those three at the top.

Whether that is a \emph{prospective} test depends on an ordering the repository cannot confirm. The channel assignment lives in the analysis script's docstring, that file was not committed before it was run, and corpus A's timestamped pre-registration does not name the channels. The strong form of this result therefore rests on our account of when two parts of one file were written. We say it here rather than only in Section 10 because this is the result the claim would inflate. What a reader can take without trusting us is weaker and still worth having: a documentary reading of ten channels is \emph{consistent with} a ranking it did not use, with no false positives among five anchors. The procedure is cheap enough for anyone to run prospectively.

A denser variant of the network, built at the source's low-confidence floor rather than at the declared cutoff, reproduces every figure: 5.1 million dyads over twelve channels, contrast 0.4132 at two to three members, 0.0000 above eighty, residual 0.4020.

\begin{figure}[!htbp]
\centering
\includegraphics[width=\textwidth]{figures/figure-7.pdf}
\caption{\textbf{The same shape in two domains, read at two levels, and a reversal when the relation saturates.} Clique contrast against group size in six bins, at zero lead time in both corpora because corpus B has no time axis.}
\end{figure}

\subsection{Task Results}

For completeness we report the accession forecasting task, since the diagnostic findings above are about instruments and a reader will reasonably ask whether the underlying prediction problem is tractable at all. A second task, the termination of a group, is in Appendix A. Table 9 gives the accession task, a dyad-level discrete-time hazard over 2,197,744 candidate periods with 6,501 event dyads, evaluated at four rolling-origins and four lead times with a cluster bootstrap over families. The unit is the dyad rather than the event: 6,597 accession events fall inside the window and 96 of them are repeats on 94 dyads that had left and rejoined, which the risk-set construction collapses onto the first because a dyad exits the risk set when it accedes. Section 5.5's flow gives that filter with the others.

\begin{table}[!htbp]
\centering
\small
\begin{tabular}{@{}
  >{\raggedright\arraybackslash}p{(\linewidth - 8\tabcolsep) * \real{0.2553}}
  >{\centering\arraybackslash}p{(\linewidth - 8\tabcolsep) * \real{0.1841}}
  >{\centering\arraybackslash}p{(\linewidth - 8\tabcolsep) * \real{0.1529}}
  >{\centering\arraybackslash}p{(\linewidth - 8\tabcolsep) * \real{0.1608}}
  >{\centering\arraybackslash}p{(\linewidth - 8\tabcolsep) * \real{0.2469}}@{}}
\toprule
\begin{minipage}[b]{\linewidth}\raggedright
\textbf{Rung}
\end{minipage} & \begin{minipage}[b]{\linewidth}\centering
\textbf{Features}
\end{minipage} & \begin{minipage}[b]{\linewidth}\centering
\textbf{Pooled AUC}
\end{minipage} & \begin{minipage}[b]{\linewidth}\centering
\textbf{Within group AUC}
\end{minipage} & \begin{minipage}[b]{\linewidth}\centering
\textbf{Calibration slope}
\end{minipage} \\
\midrule
1, duration only & 1 & 0.476 to 0.715 & 0.461 to 0.524 & -0.13 to 1.03 \\
2, plus node covariates & 22 & 0.794 to 0.911 & 0.527 to 0.632 & 0.67 to 1.02 \\
3, plus relational densities & 34 & 0.798 to 0.915 & 0.566 to 0.688 & 0.65 to 0.96, mean 0.83 \\
\bottomrule
\end{tabular}
\caption{Accession task, corpus A. Each rung must beat the one below on both metrics.}\label{tab:9}
\end{table}

The two discrimination metrics differ by about 0.25 and that gap is a finding rather than a nuisance. Pooled, the model scores 0.798 to 0.915. Holding the group fixed and asking only which candidate accedes, it scores 0.566 to 0.688. Most of the pooled discrimination is therefore between-group and between entity structure: which groups grow, and which entities join things in general. Both are legitimate questions and they are different questions, so reporting one number would obscure which was answered. This generalises past our corpus and Section 11 carries it as a recommendation: \textbf{any group membership task should report both a pooled and a within-group discrimination}. A pooled figure on a task of this shape is dominated by which groups grow, so a reader who wants to know whether the model picks the right candidate is reading the wrong number. Relational densities add 0.0089 pooled and 0.0430 within-group, nearly five times more once the group is held fixed, which is the sense in which relational structure answers a different question from node covariates. Both figures are means of the rung 3 minus rung 2 difference over the sixteen origin and lead time cells, not differences of the reported ranges. The ranges in Table 9 are minima and maxima over those cells and their endpoints differ by about 0.004, so reading the gain off the endpoints gives a different and wrong number. Under an alternative duration clock that measures eligibility from group founding rather than from entry into the observed risk set, the within-group gain is 0.0518 and the ordering of rungs is unchanged.

Calibration is reported as a number rather than a caveat: rung 3's slope averages 0.83 and never exceeds 1.0, so the model is systematically overconfident by roughly 17 percent on the logit scale.

A second forecasting task, the termination of a group, is reported in Appendix A. It is a discrete-time competing-risks hazard with interval-censored event times over 527 groups and 150 events. The 527 are the groups that contribute a period to the likelihood, against the 528 whose observation interval overlaps the window, and Section 5.5's flow separates the two. One figure from it belongs here because the diagnostic argument uses it: cause-specific discrimination is 0.628 and 0.559 against the accession task's pooled 0.798 to 0.915. Termination is therefore the harder problem on the comparable metric, and the tractability question a reader asks about the accession task does not transfer to it. The rest of that task, including its coefficient table and the two sign reversals a pooled model would average away, bears on the substantive problem rather than on the instrument, and it is in the appendix for that reason.

Table 10 and Figure 8 give the architectural comparison, four encoders over the same risk set with parameter-matched graph-free controls, at \textbf{ten} seeds each and at four. Both counts are reported because Section 9's fourth question is how much a number moves when nothing substantive changes, this comparison is the example that question was drawn from, and stopping at four seeds was the same mistake one step further out. One verdict changes.

\begin{table}[!htbp]
\centering
\small
\begin{tabular}{@{}
  >{\raggedright\arraybackslash}p{(\linewidth - 10\tabcolsep) * \real{0.1455}}
  >{\centering\arraybackslash}p{(\linewidth - 10\tabcolsep) * \real{0.1329}}
  >{\centering\arraybackslash}p{(\linewidth - 10\tabcolsep) * \real{0.2068}}
  >{\centering\arraybackslash}p{(\linewidth - 10\tabcolsep) * \real{0.1435}}
  >{\centering\arraybackslash}p{(\linewidth - 10\tabcolsep) * \real{0.2173}}
  >{\centering\arraybackslash}p{(\linewidth - 10\tabcolsep) * \real{0.1540}}@{}}
\toprule
\begin{minipage}[b]{\linewidth}\raggedright
\textbf{Encoder}
\end{minipage} & \begin{minipage}[b]{\linewidth}\centering
\textbf{Parameters}
\end{minipage} & \begin{minipage}[b]{\linewidth}\centering
\textbf{Mean AUC, 4 seeds}
\end{minipage} & \begin{minipage}[b]{\linewidth}\centering
\textbf{SD, 4 seeds}
\end{minipage} & \begin{minipage}[b]{\linewidth}\centering
\textbf{Mean AUC, 10 seeds}
\end{minipage} & \begin{minipage}[b]{\linewidth}\centering
\textbf{SD, 10 seeds}
\end{minipage} \\
\midrule
Graph free MLP, width 64 & 54,849 & 0.8602 & 0.0023 & 0.8623 & 0.0042 \\
Graph free MLP, width 96 & 80,993 & 0.8623 & 0.0064 & 0.8644 & 0.0056 \\
Graph free MLP, width 128 & 109,185 & 0.8616 & 0.0020 & 0.8630 & 0.0033 \\
Homogeneous message passing & 59,649 & 0.8593 & 0.0029 & 0.8532 & 0.0066 \\
Relation specific message passing & 78,849 & 0.8744 & 0.0065 & 0.8743 & 0.0043 \\
Temporal encoder over the graph history & 103,937 & 0.8830 & 0.0019 & 0.8822 & 0.0031 \\
\bottomrule
\end{tabular}
\caption{Architectural comparison, corpus A. Mean pooled AUC and the seed standard deviation, at four seeds and at ten.}\label{tab:10}
\end{table}

\textbf{Two margins survive and one changes sign in significance.} The standard error of a margin here is the pooled standard error of the difference of two means over \(n\) seeds each, \(\sqrt{s_{1}^{2}/n + s_{2}^{2}/n}\), with \(s\) the seed standard deviation reported in Table 10. Every multiple below therefore follows from that table and one line of arithmetic.

The relation-specific encoder beats its parameter-matched control by 0.0100 at ten seeds against 0.0121 at four. The margin shrinks and its significance rises, from 2.7 standard errors to 4.5, because the standard error halves; it wins 95 of the 120 paired origin-by-lead-time-by-seed cells. The temporal encoder beats its parameter-matched control by 0.0192 at ten seeds against 0.0214 at four, 13.4 standard errors, winning 115 of 120 cells. Widening the graph-free control buys 0.0014 from width 64 to 128 at four seeds and 0.0007 at ten, and is non-monotone in width at both counts, so neither gain is capacity.

\textbf{The homogeneous encoder is worse than its control.} At four seeds it did not look it. At four seeds its margin over the width-matched graph-free rung was \(- 0.0009\), half a standard error, winning exactly 24 of 48 cells, which reads as indistinguishable from a multilayer perceptron rather than worse than one. At ten seeds the margin is \(- 0.0091\), 3.7 standard errors, and it wins 25 of 120 cells. Homogeneous message passing over this graph is not merely uninformative, it costs accuracy against a graph-free model of the same size, at every lead time examined: \(- 0.0067\) at one period, \(- 0.0083\) at three, \(- 0.0101\) at five and \(- 0.0111\) at ten. Four seeds got the sign right and the readability of the sign wrong, which is Section 9's fourth question applied to this paper's most quoted number, and the reason both seed counts are printed above.

The relation-specific encoder's own history is the same lesson at a smaller scale. Two seeds put its margin at 0.0009, four at 0.0121, ten at 0.0100; it has the widest seed variance in the four-seed comparison, 0.0065 against the temporal encoder's 0.0019, and a margin of that size cannot be resolved by two draws. We report every parameter-matched margin against its seed spread and its cell count for this reason.

The graph's contribution is largest at the shortest lead time and erodes with distance: over the width-64 rung the temporal encoder gains 0.0227 at one period against 0.0202 at ten, and the relation-specific one 0.0178 against 0.0095, both at ten seeds. This is the reverse of the rationale usually offered for relational features at long horizons, namely that structure should matter more when covariates go stale. The erosion is monotone for the relation-specific encoder and not for the temporal one, whose profile is flat within its seed spread past the first period, so we state the direction and not a rate.

\begin{figure}[!htbp]
\centering
\includegraphics[width=\textwidth]{figures/figure-8.pdf}
\caption{\textbf{Message passing against a capacity-matched control, with the seed spread that settled it.} Corpus A accession task, mean pooled AUC over ten seeds against parameter count, six encoders, with the seed standard deviation as the error bar.}
\end{figure}

\section{Ten Ways a Leakage Diagnostic Misreports}

\textbf{The count, stated before the list.} Eleven items are lettered below and ten are demonstrated by a measurement named in the item. Item (e) is a concern that is sound a priori, that we tested, and that we could not demonstrate. It is set apart in Section 7.1 rather than counted, because a tally of the ways an instrument fails should not be padded by a way it did not. The claim is therefore ten mechanisms and one unsupported concern.

Each item states a mechanism, the measurement establishing it, and whether the second domain showed it to be general. The ten are not the same kind of thing and Table 11 says which kind each is, because a single list flattens differences that matter to a reader deciding what to carry away. Three follow from definitions or from combinatorics and hold wherever the statistic is used, items (a), (c) and (k). Three are empirical findings we take to be general, items (f), (g) and (i). Three are empirical and specific to the coalition corpus in magnitude, items (b), (d) and (h). One is a design principle, item (j).

\textbf{(a) The suite needs an operation its specification lacks.} Focal exclusion stops the group's own edges standing in for the label. Detecting back-dating requires focal isolation of those same edges, since that is the only place back-dating lives. Two operations, opposite purposes, and a specification that names only the first cannot measure the quantity it exists to measure. Established in Section 6.4. \emph{General}, because it follows from the definitions rather than from any corpus.

\textbf{(b) A saturating target relation hides the trace.} Co-membership has density 0.9954, so almost every pair of states shares something. Of the candidate-to-member dyads the pooled statistic reads as present, the focal entity is solely responsible for 0.02 percent at \(k = 1\) and 6.5 percent at the declared primary level \(k = 25\). Both are computed on this corpus rather than carried from the trade register of Section 5.3, whose own increment is a different 6.8 percent of a different relation. Thresholding does not fix it and is non-monotone in density; excluding near universal bodies leaves 0.706. \emph{Corpus specific in its magnitude}: no relation in corpus B exceeds 0.001887. \emph{Not independent of (h)}: Section 6.2 shows the clique statistic surviving strong homophily at density 0.12, failing to fire at all at 0.50 and overshooting at 0.90, so saturation is the condition under which homophily breaks it and the two items describe one interaction from two sides.

\textbf{(c) The exact clique statistic is structurally zero on large groups, at any level where it discriminates.} A ninety member group has 4,005 pairs and a clique requires all of them. On corpus A at the declared primary level \(k = 25\) the contrast runs 0.5833, 0.3617, 0.3391 and 0.2437 across the four size bins below twenty members, 0.0521 for 21 to 80, and \textbf{0.0000 above eighty}, where the positive clique rate is itself 0.0000 and 41.1 percent of the corpus's events live. \emph{General}: corpus B reproduces the same shape, 0.4124 at two to three members through 0.0000 above eighty.

\textbf{The level has to be stated.} The size dependence reverses without it. At \(k = 1\) the same statistic on the same corpus \emph{rises} with group size, from 0.0625 at two to three members to 0.4778 above eighty. Co-membership at \(k = 1\) has density 0.9954 and both classes are near-cliques, so the contrast comes from the negative rate falling with size rather than from the positive rate holding up. The mechanism in (c) is a property of the statistic where the relation discriminates, and mechanism (b) is why \(k = 1\) is not such a level for this relation. The two interact, and a claim about clique based statistics that does not name its level is false at the level a reader would assume.

\textbf{(d) Group size and sample size can trade off, and when they do no subpopulation works.} In corpus A the subpopulation where the statistic has power is also the one with 177 events, giving a standard error equal to the effect. \emph{Falsified as stated for curated group corpora in general, and restated with a condition}: corpus B has 1,399 positives at median size 3 with a contrast of 0.4124. The trade off holds when group size correlates with entity type and the small-size type is rare.

\textbf{(e) A lead-time-resolved probe with free capacity can manufacture a horizon.} Sound a priori, tested here, and not demonstrated on this corpus. Set out in Section 7.1 rather than in this list, for the reason given above.

\textbf{(f) A decaying shortcut curve is not evidence of leakage.} Predicting ten periods ahead is harder than predicting one with no leakage at all, so decay is the null. \textbf{Decay here and in Section 7.1 means one thing: the drop in the probe's area under the curve from lead time one to lead time ten, on the panel named}. Measured on corpus A's never-event panel, 6,314 positives over 311 families at one period falling to 4,043 over 233 at ten, the shortcut curve runs 0.7639 to 0.6320, monotone, a decay of 0.1319, with a mean non-linear gain of 0.0111 over the linear probe. Section 7.1's figures are two to three times larger because they are taken on the \emph{matched} panel, where negatives are drawn from the same group and period and 1,061 positives remain at one period. Matching removes the between-group signal the probe can otherwise use at every lead time, so the curve starts lower and falls further: one statistic on two panels, not two definitions. The test that separates leakage from difficulty is the exclusion contrast at each lead time, not the level of the curve. \emph{General}.

\textbf{(g) Shortcut reliance is not a property of a model.} It depends on how saturated the leaking relation is in the regime the model is evaluated in, and saturation drifts with the population: co-membership density is 0.9704 in 1948 to 1963 against 0.9954 in 1965 to 2014, with 116 organisations per year against 273. That much is measured on two large samples. The direction of the association is measured on fewer points and it is measured on a declared set. Across \textbf{fourteen} models spanning the reliance gradient, reliance and shift fragility correlate at Spearman \(- 0.667\) with \(p = 0.0092\). The model set is every cell of a grid rather than a list: the node covariate block present or absent, crossed with seven relational blocks from none through co-membership at either level to every density with the clique indicators. That matters because a rank correlation over hand-picked points can change sign if one of them moves: the same statistic over ten chosen models read \(- 0.699\) at \(p = 0.0245\). Enumerating the grid rather than adding points that help brought in the four it was missing, after which the sign held and the \(p\) value halved. We still state the negative sign as what we observed rather than as an established reversal of the usual assumption. Fourteen is not many, and the reliance spread is 0.1667 against a fragility spread of 0.0837. The claim we make from it is the weaker one the mechanism needs: reliance and fragility are not positively coupled in the way an audit that treats reliance as a model attribute assumes. \emph{General as a caution}, since the mechanism is that the interaction is between model and regime rather than a model attribute, and the caution does not depend on the sign.

\textbf{(h) The reference band cannot subtract homophily, and whether it can depends on the anchors rather than on the corpus.} Established in Section 6.5 and Table 7. Corpus A's anchors decay like the target relation, the residual falls below the threshold at a lead time whose bootstrap interval spans the whole examined range, and focal isolation on the same events says one period at most. \emph{Corpus specific as a magnitude}: on corpus B the residual is 0.4075 and the anchors are flat. \textbf{The explanation runs opposite to the intuition}. Section 6.2 produces both regimes at a known width. With anchors from the same latent space the band subtracts the homophily and the rule is exact at every homophily setting. With anchors from an independent space the rule degrades to a slope of zero on the true width. Shared structure is what makes the band work. What corpus A actually exhibits is the conjunction of shared structure with a saturated target relation, which is item (b), and it is that conjunction rather than the sharing alone that leaves a residual the rule reads as a horizon. Item (k) is the third thing the band cannot subtract, and unlike homophily it is administrative: a trace on another group's record enters the pooled contrast and is not homophily, so a band that succeeded perfectly at removing homophily would still leave it.

\textbf{(i) The trace confound grows with lead time, and it grows toward confirmation.} Prior membership spells put an entity on the record at long lead times for reasons that are not anticipation, and reaching further back is more likely to land inside an earlier spell. Presence among recurrent accessions runs 0.185 at one period to 0.330 at eight against 0.0009 among first-time ones. Unconditioned, the affected stratum shows a textbook four-period decay that is almost entirely this. \emph{General}, and the correction is to condition on first entry and report the recurrence rate.

\textbf{(j) A diagnostic whose answer is known in advance finds defects that self-consistent validation cannot.} Focal isolation asks whether a member is on the group's record before it joined, which is a question with a known answer. At its first corpus-wide application it located a coding convention defect: 149 of 373 events, 40.0 percent, were artefacts of a register that codes a member once per treaty phase, so a founding member with a later phase entry acquired an accession it never had. The defect had survived a validation suite, two independent count audits and a consistency checker. \emph{General as a design principle}, and Section 9 develops it.

\textbf{(k) The trace of an event can live on a register the audit does not open.} Focal isolation reads a group's own record, and the operation's freedom from homophily depends on that record holding the whole trace. Section 6.3 makes a curatorial wave a controlled parameter and finds the dependence is on the corpus's dating conventions rather than on the operation. Under one convention the wave adds nothing a focal reading misses, at every wave share tested. Under mixed conventions, where a neighbouring group's register is coarser, the administrative reach exceeds what the focal group's record shows, and the remedy is procedural: reading every register that feeds \(r^{\star}\) rather than the focal group's alone recovers most of the difference. \emph{General as a condition}, since it follows from the registers rather than from either corpus, and it is the term \(\lambda\) in Section 3.4's decomposition. It is also a correction to this paper's own Section 3.3, which claimed the focal record is the only place back-dating can live.

\subsection{Assessed and Not Supported: a Probe with Free Capacity}

We include this because we proposed it, tested it, and could not demonstrate it, and a list of ways an instrument fails is more informative if the entry that did not survive is visible rather than absorbed.

The positive count necessarily falls with lead time, 1,061 at one period against 508 at ten on the matched panel, so a high capacity probe has more room to overfit where the sample is smallest and can yield a decaying curve whatever the leakage is. That concern is sound a priori and we still hold it. Measured on the matched panel it is not what drives the decay here. Decay is the same quantity mechanism (f) defines, the drop in area under the curve from lead time one to lead time ten. On this panel it runs 0.2160, 0.2596 and 0.2686 across 30, 100 and 300 trees, a ratio of 1.24 between the extremes, while the \textbf{linear} probe decays 0.2597, which sits between them. A decay present at every capacity including linear is not a capacity artefact. These numbers are larger than mechanism (f)'s 0.1319 because the panel is matched and mechanism (f)'s is not, which is the comparison that paragraph now states rather than leaving to arithmetic.

Reporting a linear probe beside a non-linear one is what settles the question, and it settles it against the mechanism. The figures the claim originally rested on, a decay of 0.058 against 0.143 with a linear probe at 0.041, came from a panel whose negatives were drawn across group periods rather than within them, and that panel is void. The capacity comparison was internal to one panel, so the mechanism was unaffected by the defect and its numbers were not, and on the corrected panel they do not support the claim. The recommendation survives either way: report a linear probe alongside a non-linear one, because that is what makes the question answerable at all.

Table 11 summarises which mechanisms transfer, and what kind of claim each is.

\begin{table}[!htbp]
\centering
\small
\begin{tabular}{@{}
  >{\raggedright\arraybackslash}p{(\linewidth - 6\tabcolsep) * \real{0.0853}}
  >{\raggedright\arraybackslash}p{(\linewidth - 6\tabcolsep) * \real{0.3000}}
  >{\raggedright\arraybackslash}p{(\linewidth - 6\tabcolsep) * \real{0.2743}}
  >{\raggedright\arraybackslash}p{(\linewidth - 6\tabcolsep) * \real{0.3403}}@{}}
\toprule
\begin{minipage}[b]{\linewidth}\raggedright
\end{minipage} & \begin{minipage}[b]{\linewidth}\raggedright
\textbf{Mechanism}
\end{minipage} & \begin{minipage}[b]{\linewidth}\raggedright
\textbf{Type}
\end{minipage} & \begin{minipage}[b]{\linewidth}\raggedright
\textbf{Status after replication}
\end{minipage} \\
\midrule
a & Missing focal isolation & definitional & general, follows from definitions \\
b & Saturating target relation & empirical, corpus specific & corpus A specific in magnitude; interacts with h \\
c & Clique statistic zero on large groups & combinatorial & general, whole curve traced \\
d & Size and sample trade off & empirical, corpus specific & falsified as stated, restated with a condition \\
e & Free capacity manufactures a horizon & a priori, tested & \textbf{not demonstrated} on this corpus; Section 7.1 \\
f & Decay is not evidence & empirical, general & general \\
g & Reliance is not a model property & empirical, general as a caution & general as a caution; the sign rests on \(n = 14\) \\
h & Band cannot subtract homophily & empirical, corpus specific & corpus A specific; explanation reversed in 6.2 \\
i & Recurrence confound grows with lead time & empirical, general & general; the phenomenon is known for link prediction \citep{poursafaei} \\
j & Known answer diagnostics find defects & design principle & general design principle \\
k & Trace on a register the audit does not open & definitional, measured & general as a condition; corrects our Section 3.3 \\
\bottomrule
\end{tabular}
\caption{Generality of each mechanism, with the type of claim each makes. The type column exists because the items are not comparable: a consequence of the definitions and an empirical finding on one corpus are different objects, and a reader deciding what to carry into another domain needs to know which is which.}\label{tab:11}
\end{table}

\section{What to Do When the Register Codes One Date}

This is the practical gap the paper's own findings open, and closing it is not optional. Section 3.4's condition says focal isolation identifies a width only where a source dates the relation and the event from separate fields. Most registers do not. Corpus A does not, for 94 percent of its events. So a practitioner who reads Section 11's first instruction, measure back-dating on a group's own record, and turns to a single-dated register will find the measurement returns zero by entailment and tells them nothing. Five things can be done instead, in descending order of what they establish. The fifth is not an alternative to the first four but a condition on all of them, and Section 6.3 is what put it on this list.

\textbf{Match a second date from outside the register.} The width is identified as soon as any second dating of the same event exists, and it need not come from the same source. Section 5.3 is an instance: the trade agreement register happens to code both dates, but the same structure is available whenever an external source records when an event was \emph{decided} against when the register records it as \emph{effective}. The procedure is to join on the event key, keep only events matched in both, report the match rate, and run focal isolation against the externally dated event. The match rate bounds what the measurement covers, and an unmatched majority is a reason to report the width for a subpopulation rather than to abandon it.

\textbf{Audit a sample of events against the primary source.} Draw \(k\) events at random, return to the document the register was coded from, and record the date the document supports. This is expensive per event and cheap in total: even \(k = 50\) bounds the share of events whose coded date is displaced. Section 6.2 gives the mapping from that share to the width focal isolation would read, a displaced share of a tenth producing a presence of 0.0667 at one period. The audit is also the only procedure on this list that can detect a systematic displacement in the register as a whole.

\textbf{Report the sensitivity of the result to the horizon rather than the horizon.} If the width cannot be measured, the quantity a practitioner actually needs is still available: how much does the conclusion depend on it? Fit at each lead time under the assumption that the horizon is \(\beta\), for \(\beta\) across a declared range, and report the range of the performance figure rather than one number. A result that survives every plausible \(\beta\) needs no horizon estimate; a result that does not is a result that was resting on one. This inverts the problem in the direction the evidence allows. It establishes least about the width of the four and, on a register we could not date twice, it is nonetheless the first thing we would report.

\textbf{Read the coding rules and declare the channels.} Section 6.6 is the evidence that this works: a written declaration of what each of ten channels records placed the three contaminated channels first, second and third on clique contrast, with no false positives among five anchors. Reading is not a substitute for measuring a width, but it identifies which relations can carry the event at all, and that is the decision most exposed to being made by a statistic that cannot make it. Whether that reading was prospective is a question about our own record-keeping rather than about the procedure, and Section 6.6 answers it against us.

\textbf{Measure on every register that feeds the target relation.} Not on the focal group's alone. This is the fifth item and it applies to the four above rather than competing with them. A second date, an audited sample or a declared reading covers only the register it was taken on, and if the corpus contains more than one dating convention the widest one governs the target relation. Section 6.3 measures the cost of the narrower reading on a generator where the truth is set, and it is a mean of 0.82 to 2.68 periods and a maximum of nine. The check that says whether this matters for a given corpus is cheap and is not a statistic: count the dating conventions its sources declare. Corpus A declares two, one exact and one accurate to five-year blocks, so it is in the regime where the narrower reading undershoots. Then state the range of lead times examined and make it at least as wide as the widest convention, because Section 6.3 shows that the instrument censors at the top of its range exactly as the pooled statistics do.

What none of these does is recover a horizon from the target relation's pooled density. Section 6.2 measures why: once homophily is strong enough, the pooled statistic's output is independent of the true width.

\section{Six Questions with Externally Checkable Answers}

The largest single source of corrections in this work was not a leakage measurement. It was an audit of whether each reported number could be regenerated, and it belongs in a paper about measurement validity for a reason that is easy to state and easy to dismiss. A diagnostic's output is only as good as the artefact chain behind it, and a chain that agrees with itself is not evidence about that chain. The audit found nine defects in a single working day. \textbf{Four of them changed a leakage number}, which is the connection, and they are worth naming. A register's phase-row confusion inflated 162 group periods' size by up to a factor of two, and group size is the variable mechanism (c) is about. A shared output directory silently added four features to one reader's input, two of them built on a deliberately back-dated relation, which is leakage introduced by a filename pattern. A construction check was reported as a measurement, which is the entailment-versus-measurement error Section 3.4 now exists to prevent, committed before we named it. And a default parameter list that omitted one setting produced ``five of five'' where the calibration is six of six, which is the single number Table 3 rests on. The other five changed task results or pipeline behaviour. Appendix C lists all nine with what each was and what found it.

\textbf{Why the checks in place did not find them.} The validation suite, two independent count audits and a document-to-artefact checker all test \textbf{internal agreement}. A corpus can agree with itself while counting rows as members. A document can agree with an artefact that no script produces. A test can agree with a closed form for a code path nothing calls. What found all nine were questions whose answers lie outside the system, and the reporting standard this section is an instance of is REFORMS \citep{reforms}, whose requirement that every reported figure be traceable to the code that produced it is what the stamping below implements.

\begin{quote}
\textbf{Six questions with checkable answers.} The first four found the nine defects of Appendix C. The fifth and sixth were found by preparing this revision, and the sixth cannot be automated at all.
\end{quote}

\begin{enumerate}
\def\labelenumi{\arabic{enumi}.}
\tightlist
\item
  \textbf{What command produced this artefact?} Recorded in the artefact, not reconstructable from its filename.
\item
  \textbf{What is this derived object derived from?} A cache or panel must carry the identity of its inputs.
\item
  \textbf{Does anything actually call this?} A tested capability is not a used one.
\item
  \textbf{How much does this number move when nothing substantive changes?} Two seeds put an architectural margin at 0.0009, four at 0.0121 and ten at 0.0100, and four made an encoder that is significantly worse than its control look like a tie.
\item
  \textbf{Does any second code path compute this number, and does it agree?} The only question on the list that requires writing something twice, and the one that caught a density wrong by exactly a factor of two.
\item
  \textbf{Does the sentence around this number attribute it to the right object?} No check we can write answers this. It needs a reader who is not us.
\end{enumerate}

We stamp every artefact with the command that wrote it and a three-integer fingerprint of the corpus it was computed on, and the consistency checker requires the stamp, so a new ad hoc artefact fails a check rather than waiting to be cited. That is the \emph{how} and \emph{where} of provenance in the database sense \citep{datasheets,modelcards} at its crudest, and it is enough for the first three questions. One boundary is worth stating: after the stamping was in place, two numbers in one section were quoted for several hours after a refactor removed the code that produced them. Provenance stamping does not prevent a claim from losing its source; it makes the loss detectable, and the check must be extended each time a claim becomes load bearing.

\textbf{The fifth question came from computing something twice by accident.} Writing the negative-sampling audit of Section 5.5 required the density of corpus B's target relation, and the figure it produced disagreed with the standing one by a ratio of exactly 2.0000. The source ships both directions of every interaction and the density block divided a directed row count by an unordered dyad count, so every corpus B density in the paper was twice its value, in all twelve channels at both thresholds. The clique statistics were untouched, because the code that computes them builds a dyad key set and collapses the duplicate directions, which is why nothing else moved and why the defect could sit in one summary table indefinitely. Hence: \textbf{does any second code path compute this number, and does it agree?} It is the argument Section 10 makes for the independent implementation of the density code, arriving from a defect rather than from a threat. It is also the one question here that requires writing something twice.

\textbf{The sixth question came from outside.} No version of it is cheap. Two external reads of this manuscript produced twelve defects between them, of a class none of the five reaches. Three were pairs of counts that differ for a correct reason nobody had written down: 528 groups against 527, 7,146 events against 7,084, 6,501 dyads against 6,597. Every one of those six numbers is right. The filters that separate them are an active membership row inside the window, the lead time the count is taken at, and one accession per dyad, and none of the three was stated. Chasing them found a fourth pair, 151 coded terminations against the 150 the panel scores. Three were terminological: an estimator called a lower bound that overshoots what it bounds, a ratio between a probability and a density contrast described as a fraction of a width, and a second domain called a replication when it has no time axis. Six were citations, including three release notes that predate the release the corpus was built from, so a reader following the citation would have found different data. Appendix B records each with what forced it.

Every one of these is a \textbf{claim-to-artefact mismatch rather than an artefact defect}. The checker tests whether each registered value appears somewhere in the manuscript, that every value anchored to a unit matches everywhere it appears, and that every section and table a sentence points at exists. It cannot test, and cannot be made to test, whether the sentence around a value attributes it to the right object. That is the sixth question, and it is the only one on the list with no cheap version. The case for it is now quantitative: two outside reads produced twelve defects of this class and led us to a thirteenth, and the internal checks that found the other nine found none of them.

\section{Threats to Validity}

\textbf{A single implementation, cross-checked against a second.} What the cross-check found was a specification gap rather than an arithmetic error. The clique contrast and group density were computed by one implementation across all four testbeds, which is what makes cross domain comparison interpretable and also meant an implementation error would have been common to all results. We have since written a second implementation from Section 3's definitions, sharing nothing with the first. It stores dyads in Python dictionaries where the original packs them into sorted integer arrays and answers membership by binary search. And it enumerates every pair of every candidate group where the original decomposes the group into a cached member-only part and a per-candidate star, which is the optimisation most able to be subtly wrong. The two agree \textbf{exactly} on every value in every testbed: 74,425 densities on the synthetic corpus, 784,585 star densities behind Table 7, 39,582 candidate-group densities behind Table 8, and 101,600 densities on corpus B, where a third implementation from a separate code base also agrees exactly. Table 3's calibration was re-derived end to end through the second implementation and recovers the known width in 6 of 6 settings, with an estimate identical to the original's at every setting.

What the cross-check bought was not an arithmetic correction. It was the discovery that one convention had never been written down clearly enough to be re-implemented. The second implementation read ``level 1'' literally, as weight at least one, which empties the voting-agreement relation and discards half of trade flows, and it disagreed with the original by exactly 1.0 on those two stars. The original was right and the prose was silent; Section 4 now states the convention. We report this because it is the class of defect a self-consistent validation suite cannot find and Section 9's first four questions do not reach either: both implementations were written from the same prose, and the disagreement located an ambiguity in the prose.

\textbf{The interval is withdrawn.} Its upper end is not identified. Section 6.1 gives the two failure modes: as specified it is unreachable in half the synthetic settings and never fires on corpus A, and repaired as a band crossing it fires everywhere but coincides with \({\widehat{\Delta}}^{\star}\) in every synthetic cell and falls \emph{below} it on corpus A. The two statistics do not order as the bracket requires. This is a property of the definition rather than of a corpus, and predictions bounded above cannot be tested against either version.

\textbf{What survives is not a bound.} Section 3.2 sets this out and it belongs here too. The number the suite reports errs upward, in a third of resamples at the widest synthetic setting, so it does not bound \(\Delta^{\star}\) from below. It is conservative under the asymmetric loss Section 3.2 names, and it is not conservative under a symmetric one. A practitioner who needs a guaranteed bound in either direction does not have one here, and we know of no statistic in this family that would supply one.

\textbf{Focal isolation is complete only under one dating convention.} It is unbiased in homophily everywhere. Section 6.3 measures the condition rather than assuming it. Where every register in a corpus back-dates by one law, reading the focal group's record captures the whole administrative trace at every wave share tested. Where conventions differ across groups, which is corpus A, a trace of the focal event reaches the target relation through another group's record and a focal reading does not see it; reading every register recovers most of the difference and not all of it. Two limits follow. The first is that Section 11's instruction has to name the wider reading, which it now does. The second is that we cannot verify on corpus A that the wider reading is sufficient, because there the true width is what we are trying to measure; the sufficiency is measured on the generator and inherited here.

\textbf{No comparison against another leakage diagnostic.} We report the suite against ground truth and against itself and not against a competing instrument. The reason is that we could not find one to compare against. The leakage literature supplies checklists and taxonomies rather than estimators \citep{kaufman,kapoor_leakage,reforms}, and the temporal-graph evaluation literature supplies protocols and negative samplers rather than horizon estimates \citep{ogb,tgb,poursafaei}. We are aware of no published statistic that takes a relation and a lead time and returns how far a register's dating reaches. The four estimators graded against each other in Table 4 are therefore variants of our own rule, including the one a practitioner who skipped Section 3.2 would compute, and that is the nearest thing to a baseline the field currently offers. If such a statistic exists we would want the comparison, and its absence is a limitation of this paper rather than a claim about the literature.

\textbf{One domain has no time axis.} So the second domain is a scope test rather than a replication. Corpus B tests the size and homophily mechanisms but cannot test anything about lead time, so mechanisms (e), (f) and (i) rest on corpus A alone. It follows that no central measurement of this paper has been reproduced on a second corpus: the horizon is measured in lead times and corpus B has none. Section 6.6 calls what it does a scope test for that reason, since calling it a replication would promise more. Mechanism (e) rests on an a priori argument rather than a measurement, since the measurement we had for it came from a retracted panel and the corrected panel does not reproduce it. Section 6.2's generator supplies a lead-time axis with a known truth and partly closes this gap, but it is synthetic. It can show that homophily destroys identification, and the question is whether any real corpus reaches the strength at which it does. Section 6.2 answers that on a scale both sides have: two kernels sharing no functional form break the pooled statistic at anchor clique contrasts of 0.0152 and 0.0186, and corpus A's strongest anchor is 0.0259. So the collapse regime is not past anything real. What the generator still cannot do is show that corpus A's own phantom is caused by homophily, because the trace there is the unknown, so this narrows the threat rather than removing it.

\textbf{The central field corpus cannot falsify the pooled reading.} The reason is structural. Section 6.5 sets this out and it belongs in this list. Corpus A codes one date for 94 percent of its events, so focal isolation is silent there by construction. The one stratum where it does measure lies entirely before the analysis window, so the pooled statistic cannot be read on those events at any positive lead time without reintroducing a coverage break. The consequence is that the falsification of the pooled reading against a known width rests on the trade register alone, with 759 agreements in a domain that is not the paper's primary one. The coalition corpus contributes a magnitude, an indistinguishability result and a non-identification result rather than a refutation. A register that coded two dates over a corpus of corpus A's size would settle what neither of ours can, and we know of none.

\textbf{The anticipation channel is not measured.} Focal isolation measures back-dating introduced by what a source records as a date, and it can do so only where a source codes two dates. The corpora that code one date return zero by entailment rather than by measurement, as Table 1 records row by row; Section 8 sets out what is available instead. Anticipation through channels other than the focal group's own record, an entity establishing bilateral ties before joining, remains unmeasured, and no contrast of this form separates it from homophily, which Section 6.2 supports with a measured bias rather than with reasoning alone.

\textbf{The magnitudes have intervals; the design decisions do not.} Every headline diagnostic magnitude is reported with a percentile interval from a bootstrap over the declared resampling unit. One of those intervals is wide enough to matter: the lead time at which the pre-specified rule declares the horizon closed on corpus A ranges over the whole examined interval under resampling. Two exceptions were closed in this revision and one remains. Table 4's transition rows are five seeds and its whole grid is three, under three homophily kernels; Table 10's architectural margins are ten seeds rather than four. What carries no uncertainty statement is the set of choices upstream of the statistics: the anchor set, the level ladder, the comparison strata, the window. Those were fixed in writing beforehand, which prevents them from being tuned to the result and does not make them right. A different documentary reading of the same five relations would produce a different band, and we have no way to put an interval on that.

\textbf{The declaration is timestamped for one corpus and not for the other.} Section 5.1 sets this out and it belongs in this list. Corpus A's pre-registration is in version control with dates that precede the data, and the two predictions it lost are in it. Corpus B's channel declaration is not, so Section 6.6's prospective reading is a claim about our own procedure that a reader cannot check, and we have downgraded what we claim from it accordingly. This is the class of defect Section 9's fifth and sixth questions describe, found on the paper's own most quotable result, and it was found by asking what the artefact behind a claim actually supports rather than by any check we had running.

\textbf{No external replication of the code or the corpus.} The nine defects of Section 9's main audit, the density double count and the second implementation of this section's first paragraph were all produced by the authors working on their own code. The boundary noted in Section 9 is an instance of the resulting ceiling. The second implementation narrows the exposure for one component and does not remove it: it was written by the same people from the same prose, which is why the gap it found was a gap in the prose. The one exception is the twelve claim-to-artefact mismatches Section 9 reports last, which came from an external read of the manuscript. Those establish that an outside reader finds a class of defect we cannot, and they are a read of the paper rather than a replication of the work. No one outside the author group has run the pipeline or rebuilt the corpus.

\section{Conclusion}

Relational features for group membership prediction are often drawn from the register that supplies the labels, and registers date their contents by conventions that can place an edge before the event it belongs to. We formalised the reach of that contamination as a leakage horizon and built a diagnostic suite to estimate it. The suite reports one number rather than an interval: the upper end we specified is not identified, either unreachable as written or, repaired, indistinguishable from the lower end on synthetic data and below it on our field corpus, and we report that as a result about the definition. Against a width set by a generator that number recovers it in 6 of 6 settings, in 67 to 100 percent of bootstrap replicates depending on the setting, and it never undershoots. That is why we call it a conservative estimate rather than a lower bound: an estimator that exceeds the quantity it was named for is safe only under a stated asymmetry of loss, and we state it. Against a register that codes the width itself, the back-dated dyads are 6.8 percent of the relation and the pooled statistic the suite specifies is dominated by the other 93.2 percent.

The reason is that a density over the whole target relation cannot distinguish a trace of the event from a persistent similarity between the acceding entity and the members. That is a measurement rather than an argument. On a generator where the width is set and homophily is a dial, the pooled statistic's estimate loses all dependence on the true width once homophily is strong enough, its least-squares slope on that width falling from 1.00 to 0.00. Focal isolation on the same corpora stays exactly unbiased. The anchor band repairs the pooled statistic when the anchors express the same latent structure as the target relation and fails when they do not, which is the opposite of the direction the mechanism is usually stated in. On our coalition corpus the specified diagnostic reports a contrast still at 0.1490 after eight years and indistinguishable there from a relation that cannot record the event, while the group's own record carries nothing past one year. The lead time at which the pre-specified rule declares the horizon closed is five on the point estimate and ranges over the whole examined interval under resampling, so the phantom is real in most samples and its width is not identified.

The practical guidance is specific, and it now includes the two cases earlier versions of this paper left open.

\begin{itemize}
\tightlist
\item
  \textbf{Measure back-dating on a group's own record, not on a pooled density, and on every register that feeds the target relation rather than on the focal group's alone.} Section 6.3 measures what the narrower reading costs where dating conventions differ across groups.
\item
  \textbf{Where the register codes one date}, which is the common case, Section 8 gives five things to do instead, of which reporting the sensitivity of the result to an assumed horizon is the one we would do first.
\item
  \textbf{Condition any trace statistic on first entry} and report the corpus's recurrence rate beside any horizon.
\item
  \textbf{Report a linear probe alongside any non-linear one}, which is what showed us that probe capacity was not driving the decay we had attributed to it, and report an exclusion contrast at each lead time rather than the level of a decay curve. Say which panel a decay is taken on, because a matched panel and an unmatched one give figures that differ by a factor of two.
\item
  \textbf{Stratify any clique based statistic by group size}, state the level it is read at, and state the share of the corpus at zero power.
\item
  \textbf{Declare the target relation and the anchors from what each source records}, before computing anything, because that reading identified all three leaking channels of a ten-channel network with no false positives among five anchors. Then deposit the declaration somewhere that dates it. The reading is the cheap half and the dating is the half we got wrong.
\item
  \textbf{Put an interval on every magnitude, including the lead time a rule returns}: ours was the number that changed most under resampling and it was the number we had reported without one. Report the censoring rate beside any slope taken over a bounded range, because a slope of zero and a range too narrow to hold the estimator look identical otherwise.
\item
  \textbf{Report both a pooled and a within-group discrimination for any group membership task.} Ours differ by about 0.25, and either number alone answers a different question from the one a reader will assume.
\item
  \textbf{Report every parameter-matched architectural margin against its seed spread}, and use enough seeds that the spread is estimable. Two seeds put one of our margins at 0.0009, four at 0.0121 and ten at 0.0100, and four also made an encoder that is significantly worse than its control look like a tie. Four is not enough.
\end{itemize}

Eight of our own claims did not survive the tests we ran against them, and three of those tests existed because a reader outside the author group asked for them. Appendix B lists all twenty changes this paper's own tests forced, with what each was and what forced it, so that the body can state conclusions and a reader who wants the audit trail can have it in one place. We regard all eight as the tests working. Take a set of claims about curated group corpora and put it through four things: an adversarial second domain, a generator built to reproduce its own failure mode, a seed count large enough to see its own noise, and a reader whose job is to disbelieve it. A set that survives all four is a set that was not adequately tested.

\section{Data and Code Availability}

\textbf{Code.} The corpus construction, the diagnostic suite, the survival estimators, the synthetic generators and every analysis script are released as one repository at \url{https://github.com/McahitKutsal/Leakage-Horizons}, under the MIT licence. All reported figures are regenerated by named scripts from the stored results, and every stored result carries the command that produced it together with a three-integer fingerprint of the corpus version it was computed on. A consistency checker requires that stamp and compares every registered number, every unit-anchored quantity and every section and table reference in the manuscript against the artefacts. It is released with the rest, so that a reader can run the check rather than take its output on trust. The second, independent implementation of the density code reported in Section 10 is released alongside the first, together with the cross-check that compares them, because a claim that two implementations agree is only checkable if both are available.

\textbf{The declaration.} The artefact Section 5.1's protocol rests on is deposited separately at \url{https://doi.org/10.5281/zenodo.22110800}, with a checksum for each file, and it contains two documents whose evidential status differs. \texttt{pre-registration.md} is corpus A's declaration as the commit that froze it left it; the repository's history carries the three commits that wrote and amended it, all of them before any data was materialised, and the manifest records their hashes and dates. \texttt{corpus-b-channel-declaration.md} is the assignment of the ten interaction channels to target, contaminated and anchor, and it is marked in the artefact itself as \textbf{not timestamped}. The file that carried it was not committed before it was run, and the pre-registration does not name the channels. Section 6.6 states the consequence where the result appears. Depositing the artefact now fixes its text and does not make it evidence about the order in which it and the statistic were produced, and we would rather say so than let a deposit date imply otherwise.

\textbf{Data.} No third-party data is redistributed. Every source is public and is obtained by a fetch script in the repository which records the release it retrieved, and none of them required an institutional affiliation of the authors. Six sources are used and each is cited in the reference list with the release the corpus was built from: the Correlates of War IGO membership register, version 3; the ATOP alliance register, version 5.1, member-level file; the DESTA trade agreement register; the Militarised Interstate Disputes register, MID5; the STRING interaction network, v12.0, \emph{Homo sapiens}; and the CORUM and Complex Portal complex catalogues, at CORUM release 5.3 and the Complex Portal current human release, both retrieved on 2 August 2026. Ideal-point estimates for United Nations voting are obtained from the replication archive of the article cited for them. The derived corpus is not redistributed because it is reconstructible from those sources by the scripts, and the repository records the fingerprint of the version every number was computed against so that a rebuild can be compared rather than assumed identical.

\textbf{The three bracketed identifiers above are assigned at release and deposit and will be inserted at acceptance. Every other statement in this section describes what the artefacts already contain rather than what we intend.}

\section{Declaration of Competing Interest}

The authors declare no competing interests.

\section{CRediT Authorship Contribution Statement}

Section 9 reports an audit of the authors' own work and Section 10 records that no external replication exists, so who performed which part of that audit is part of the claim rather than an administrative appendix. This statement is written to be read that way, and it is deliberately conservative: where a role would strengthen a claim in the paper, it is assigned only where the released repository's history supports it.

\textbf{M\"ucahit Kutsal:} Conceptualization, Methodology, Software, Formal analysis, Data curation, Validation, Writing (original draft), Writing (review and editing). The version-control history of the released repository records a single committer for the corpus construction, the diagnostic suite, the survival estimators, the synthetic generators and the analysis scripts, which is the evidence for the software, formal analysis and data curation roles. The same history is the evidence for the pre-registration dates Section 5.1 reports.

\textbf{Alpaslan Erda\u{g}:} Conceptualization, Methodology, Investigation, Writing (review and editing).

\textbf{U\u{g}ur \"Ozbek:} Methodology, Formal analysis, Writing (review and editing).

\textbf{Do\u{g}a Ko\c{c}yi\u{g}it:} Investigation, Data curation, Writing (review and editing).



\textbf{What this statement does not claim.} Every audit Section 9 reports, including the second implementation of the density code in Section 10 and the two count audits, was performed by an author who had written the code being audited. No audit in this paper was independent in the sense that matters, and Section 10's threat is stated at the level of the author group for that reason. The two external readings of the manuscript that Section 9 reports last were not performed by any author, and they are the only part of this work that had a reader outside the group; they are readings of the manuscript and not replications of the pipeline.

\acks{This work received no external funding and no grant supported it.

We thank \"Om\"ur Ko\c{c}yi\u{g}it and Ay\c{s}e Gizem Erg\"or, who contributed investigation and validation and read the manuscript.

We thank the projects that maintain the registers this paper depends on and release them publicly: the Correlates of War project for the intergovernmental organisation and dispute data, the Alliance Treaty Obligations and Provisions project, the Design of Trade Agreements project, the STRING consortium, CORUM, and the EMBL-EBI Complex Portal. Nothing in this paper would be checkable without their release practices, and the dating conventions we spend the paper measuring are documented by those projects rather than discovered by us.

Two readings of this manuscript by readers outside the author group produced eleven of the corrections we report, including the objection that became Section 6.3 and the terminology of the horizon itself. We are grateful for them, and Section 10 records that they remain the only part of this work reviewed from outside.

\emph{Disclaimer.} The views expressed in this paper are the authors' own and do not represent the position of the Republic of T\"urkiye Ministry of Foreign Affairs with which the authors are affiliated.}

\newpage

\appendix
%% The appendix tables are cited as A.1, B.1 and C.1. Without this they number 12, 13 and 14.
\numberwithin{table}{section}

\section{The Termination Task}
\label{app:A}

Reported here rather than in Section 6.7 because it concerns the substantive problem, the survival of international organisations, rather than the measuring instrument the paper is about. It is retained in full because two of its findings are methodological. The interval formulation is mandatory rather than optional, because the register's own coding rule declares an organisation defunct after ten consecutive years without a plenary session. And thirty events whose cause the source does not state enter through the likelihood's cause-unknown branch rather than being dropped or assigned a cause.

Table A.1 gives the termination task, an interval-censored competing-risks hazard over 527 groups and 16,067 group periods with 150 events. Both counts are the estimated panel's rather than the corpus's, and Section 5.5's flow gives the filters: 528 groups have an observation interval overlapping the window and 527 contribute a period to the likelihood. The one difference is a group with no active membership row inside it, and that same group carries one of the 151 coded terminations inside the window, which is why the panel scores 150.

\begin{table}[!htbp]
\centering
\small
\begin{tabular}{@{}
  >{\raggedright\arraybackslash}p{(\linewidth - 4\tabcolsep) * \real{0.3208}}
  >{\centering\arraybackslash}p{(\linewidth - 4\tabcolsep) * \real{0.2751}}
  >{\centering\arraybackslash}p{(\linewidth - 4\tabcolsep) * \real{0.2521}}@{}}
\toprule
\begin{minipage}[b]{\linewidth}\raggedright
\textbf{Covariate}
\end{minipage} & \begin{minipage}[b]{\linewidth}\centering
\textbf{Dissolution}
\end{minipage} & \begin{minipage}[b]{\linewidth}\centering
\textbf{Superseded}
\end{minipage} \\
\midrule
Log duration & +0.095 & -0.143 \\
Log size & -0.392 & -0.120 \\
Internal cohesion & +0.292 & -0.081 \\
Mean regime type & -0.093 & +0.001 \\
Regime heterogeneity & -0.109 & -0.070 \\
\bottomrule
\end{tabular}
\caption{Termination task, corpus A. Cause specific coefficients, standardised covariates, logit scale.}\label{tab:A-1}
\end{table}

Two covariates reverse sign between the two causes, log duration and internal cohesion, so a pooled model would average both reversals toward zero. Duration reads cleanly: old groups fade while young ones are replaced or merged. A reversal is counted only when both coefficients exceed 0.01 in absolute value, since a sign flip at 0.001 is a rounding artefact rather than an effect. Any cause discrimination is 0.610, cause-specific 0.628 for dissolution and 0.559 for supersession, which is the comparison Section 6.7 draws against the accession task. Thirty of the 150 events have a cause the source does not state and enter through the likelihood's cause-unknown branch rather than being dropped or assigned a cause.

\section{What Changed, and What Changed It}
\label{app:B}

Section 5.1 requires that an amendment to a declared protocol be visible rather than silent, and the body marks each one with a tag rather than narrating it. A reader meeting a result for the first time should not have to sort what is true from what we used to think. This is where the sorting is done. Twenty changes are listed. Eight are substantive enough that Section 11 counts them; the rest are corrections of naming, arithmetic or presentation. Each row says what the claim was, what it is, and what forced the change, and the fourth column is the point: not one of these was found by reading the manuscript for plausibility.

\begingroup\small
\begin{longtable}{@{}
  >{\raggedright\arraybackslash}p{(\linewidth - 6\tabcolsep) * \real{0.0663}}
  >{\raggedright\arraybackslash}p{(\linewidth - 6\tabcolsep) * \real{0.3034}}
  >{\raggedright\arraybackslash}p{(\linewidth - 6\tabcolsep) * \real{0.3232}}
  >{\raggedright\arraybackslash}p{(\linewidth - 6\tabcolsep) * \real{0.3071}}@{}}
\caption{Every change this paper's own tests and its readers forced, with what forced each. ``Own test'' means a measurement we designed and ran; ``outside read'' means a reader who is not an author.}\label{tab:B-1}\\
\toprule
\begin{minipage}[b]{\linewidth}\raggedright
\end{minipage} & \begin{minipage}[b]{\linewidth}\raggedright
\textbf{Was}
\end{minipage} & \begin{minipage}[b]{\linewidth}\raggedright
\textbf{Is}
\end{minipage} & \begin{minipage}[b]{\linewidth}\raggedright
\textbf{Forced by}
\end{minipage} \\
\midrule
\endfirsthead
\toprule
\begin{minipage}[b]{\linewidth}\raggedright
\end{minipage} & \begin{minipage}[b]{\linewidth}\raggedright
\textbf{Was}
\end{minipage} & \begin{minipage}[b]{\linewidth}\raggedright
\textbf{Is}
\end{minipage} & \begin{minipage}[b]{\linewidth}\raggedright
\textbf{Forced by}
\end{minipage} \\
\midrule
\endhead
\bottomrule
\endlastfoot
R1 & The horizon is an interval, bracketed below by the clique statistic and above by rank position with single-feature discrimination & The interval is withdrawn. The upper end is unreachable in half the synthetic settings, never fires on corpus A, and repaired as a band crossing coincides with the lower end in all 72 synthetic cells and falls below it on corpus A & Own test: the synthetic grid, Section 6.1 \\
R2 & The reported number is a lower bound on the horizon & It is a conservative estimate under a stated asymmetric loss, written \({\widehat{\Delta}}^{\star}\). Against a known width it never undershoots and exceeds the target in a third of resamples at \(\beta = 4\), so it does not bound the horizon from below & Own test, named by an outside read: Sections 3.2 and 6.1 \\
R3 & The anchor band fails because corpus A's anchors share the target relation's structure & Shared structure is what makes the band work; independent structure is what breaks it. Mechanism (h) is the conjunction of shared structure with a saturated target relation & Own test: the two anchor regimes at a known width, Section 6.2 \\
R4 & Group size and sample size trade off in any curated group corpus, so no subpopulation offers both & False as a general claim. Corpus B has 1,399 positives at median size 3 with a contrast of 0.4124. The trade off holds where group size correlates with entity type and the small-size type is rare & Own test: the second domain, Section 6.6 \\
R5 & A lead-time-resolved probe with free capacity can manufacture a horizon, and does here & Sound a priori, tested, and not demonstrated. Set apart in Section 7.1 and not counted among the mechanisms & Own test: the capacity comparison on the matched panel \\
R6 & The focal group's own record is the only place administrative back-dating can live & True of the focal event and of the target relation under one dating convention; false under mixed conventions, where the trace reaches \(r^{\star}\) through another group's record. Adds mechanism (k) and the term \(\lambda\) to Section 3.4 & Outside read, then own test: Section 6.3 \\
R7 & The suite is three components & Four. Focal isolation was absent from the specification and the design cannot do without it & Own test: Section 6.4 \\
R8 & On the trade register the pooled statistic was said to recover one tenth of the printed width & It does not estimate a width. The back-dated dyads are 6.8 percent of the relation and the pooled density is dominated by the other 93.2 percent; the quotient of 9.77 is between a probability and a dimensionless contrast & Outside read: units, Section 6.4 \\
R9 & Corpus B is a replication & A scope test. It has no time axis, so no central measurement of this paper has been reproduced on a second corpus & Outside read: Sections 6.6 and 10 \\
R10 & Corpus B's target relation has density 0.0023 and the cross-domain factor is 431 & 0.001153 and 864. A directed row count was divided by an unordered dyad count, doubling every corpus B density in all twelve channels at both thresholds & Own test: a second code path computed the same density, Section 9 \\
R11 & Raising homophily drives the pooled estimate's slope on the true width to zero & Unchanged, and now separable from its range. Every unit slope in Table 4 is uncensored; the collapsed cells are censored in all six widths, and over a range two and a half times as wide the slope goes to about \(- 1\) rather than back to 1 & Outside read, then own test: Table 5 \\
R12 & Homogeneous message passing is indistinguishable from a capacity-matched graph-free control & It is worse than one, by 0.0091 at 3.7 standard errors, losing 95 of 120 cells. Four seeds gave the sign and not its readability & Own test: ten seeds, Section 6.7 \\
R13 & Reliance and fragility correlate at \(- 0.699\), \(p = 0.0245\), over ten models & \(- 0.667\), \(p = 0.0092\), over fourteen, the models being every cell of a declared grid rather than a hand-listed set & Outside read, then own test: mechanism (g) \\
R14 & A written declaration of ten channels placed the three contaminated ones first, second and third, prospectively & The ranking stands and the prospectiveness does not. Corpus A's declaration is timestamped in version control; corpus B's channel assignment is not, so the ordering rests on our account & Own test: preparing the declaration for deposit, Sections 5.1 and 6.6 \\
R15 & The leakage horizon sits inside Kapoor and Narayanan's type L1 & Under L2, features that are not legitimate, as a subtype L2 does not name. All eight leaves are checked in Section 2 & Outside read: Section 2 \\
R16 & Six counts reported without the filters that separate them: 528 and 527 groups, 7,146 and 7,084 events, 6,597 and 6,501 dyads & All six are right and every filter is named in Section 5.5's flow. Chasing them found a fourth pair, 151 coded terminations against the 150 the panel scores & Outside read, then own test: \texttt{audit\_filter\_chain.py} \\
R17 & ``Decay'' used for two quantities that differ by a factor of two & One quantity, the drop in area under the curve from lead time one to ten, with the panel named. Mechanism (f) is on the never-event panel and Section 7.1 on the matched one & Outside read: mechanism (f) \\
R18 & Probe capacity drives the decay: 0.058 against 0.143, linear at 0.041 & Those figures came from a panel whose negatives were drawn across group periods, which is void. On the corrected panel the decay is present at every capacity including linear & Own test: the panel reconciliation \\
R19 & Table 7 reports a clique contrast & A density contrast. The caption named the wrong statistic for the right numbers & Outside read: Table 7 \\
R20 & Homophily is an exponent on latent similarity, and the collapse is reported against that exponent & One of three kernels, and the collapse is reported against the realised anchor contrast instead, because the three kernels scale the exponent differently and a threshold in it is not comparable across them & Outside read, then own test: Section 6.2 \\
\end{longtable}
\endgroup

\textbf{What this table is evidence for, and what it is not.} It is not a claim that the paper is now correct. Eleven of the twenty rows record an outside read, ten where the read came first and one where it named what our own test had found. Section 10 records that no one outside that group has run the pipeline or rebuilt either corpus, so the class of defect an external \emph{replication} would find is untouched by any of this. What the table does show is the rate: two readings of one manuscript produced twelve defects, and led us to a thirteenth, that a validation suite agreeing with itself, two count audits, a document-to-artefact checker and the five questions of Section 9 had all passed. That rate is the argument for Section 9's sixth question, and it is the reason this appendix is a table rather than a paragraph saying the paper was carefully checked.

\section{The Nine Defects of the First Audit}
\label{app:C}

Section 9 argues from the rate at which these were found rather than from their content, so the content is here. Three came from writing a script for a result that no code path reproduced, and each of those three found a further defect in the writing; six came from the same audit extended to the rest of the pipeline. The column that matters is the last: every one of them passed a validation suite that agreed with itself.

\begingroup\small
\begin{longtable}{@{}
  >{\raggedright\arraybackslash}p{(\linewidth - 4\tabcolsep) * \real{0.0691}}
  >{\raggedright\arraybackslash}p{(\linewidth - 4\tabcolsep) * \real{0.3799}}
  >{\raggedright\arraybackslash}p{(\linewidth - 4\tabcolsep) * \real{0.5510}}@{}}
\caption{The nine defects, what each was, and what it changed.}\label{tab:C-1}\\
\toprule
\begin{minipage}[b]{\linewidth}\raggedright
\end{minipage} & \begin{minipage}[b]{\linewidth}\raggedright
\textbf{Defect}
\end{minipage} & \begin{minipage}[b]{\linewidth}\raggedright
\textbf{What it changed}
\end{minipage} \\
\midrule
\endfirsthead
\toprule
\begin{minipage}[b]{\linewidth}\raggedright
\end{minipage} & \begin{minipage}[b]{\linewidth}\raggedright
\textbf{Defect}
\end{minipage} & \begin{minipage}[b]{\linewidth}\raggedright
\textbf{What it changed}
\end{minipage} \\
\midrule
\endhead
\bottomrule
\endlastfoot
1 & A register's phase-row confusion, a third instance, put 298 redundant membership rows into the corpus, 596 rows standing in duplicate pairs & Inflated 162 group periods' size by up to a factor of two. \textbf{A leakage number}: group size is what mechanism (c) is about \\
2 & A mislabelled coefficient key still present in an artefact a day after the text had been corrected & The termination coefficient table, and the fact that a corrected document does not correct its artefact \\
3 & A reversal detector scanning only the covariates, which drops the baseline term & One of the two sign reversals the termination result is scored on \\
4 & A construction check reported as though it were a measurement & \textbf{A leakage number}: the entailment-versus-measurement error Section 3.4 exists to prevent, committed before it was named \\
5 & Left truncation declared in the design, implemented in the estimator, verified against closed forms to \(10^{- 13}\), exercised by tests whose names contain the word truncated, and wired to zeros by every caller & 183 of 528 groups predate the window with a median true age of 13 periods and a maximum of 150, carrying 45.7 percent of the group periods, all on the wrong clock \\
6 & A shared output directory read by three scripts using a filename pattern rather than a declared list, which changed behaviour silently when a fourth script wrote two new files into it & \textbf{A leakage number}: a device-side assertion in one reader, four extra features in a second including two built on a deliberately back-dated relation, and a wasted read in the third \\
7 & A derived cache keyed on the analysis window rather than on the inputs it was derived from & Any result computed after a corpus rebuild, silently, until the key was changed \\
8 & A default parameter list that omitted one setting & \textbf{A leakage number}: a script reported five of five where the calibration is six of six, which is the number Table 3 rests on \\
9 & The consistency checker's entry for that same claim read an unrelated field and emitted its string regardless & The check that should have caught defect 8 \\
\end{longtable}
\endgroup

Defects 1 to 4 were found by writing the script that should have produced a stored result and comparing. Defects 5 to 9 were found by asking the first four questions of Section 9 of the rest of the pipeline. Nine defects, one working day, and the instruments designed to find defects had found none of them.

\vskip 0.2in
\bibliography{penumbra}

\end{document}
